\hypertarget{red-hot-rebis-user-guide}{%
\section{Red-Hot Rebis --- User Guide}\label{red-hot-rebis-user-guide}}

\textbf{Author:} Lando⊗⊙perator \textbf{Version:} v2.4.0 --- rebis.
Edition \textbf{Platform:} \texttt{import\ rebis} (namespace package) ·
\texttt{python3\ rebis.py\ \textless{}command\textgreater{}}
\textbf{Location:} \texttt{/home/mrnob0dy666/imsgct/red-hot\_rebis/}

\begin{center}\rule{0.5\linewidth}{0.5pt}\end{center}

\hypertarget{rebis.-namespace-organization-v2.4.0}{%
\subsection{rebis. Namespace Organization
(v2.4.0)}\label{rebis.-namespace-organization-v2.4.0}}

All tools are now consolidated under the
\texttt{rebis.\textless{}domain\textgreater{}} namespace package.

\begin{verbatim}
import rebis

# Paraconsistent kernel — genetics, hadron physics, ligand design, serpent rod
rebis.p4ra.Belnap(True, False)              # Belnap FOUR logic
rebis.p4ra.SerpentRodV2(...)                # Protein design
rebis.p4ra.encode_site_from_residues(...)   # Active site encoding
rebis.p4ra.generate_ligand_smiles(...)      # De-novo ligand generation

# Other domains
rebis.ch3mpiler.forward("CC(=O)O")          # Molecular compiler
rebis.clink.CLINK_LAYERS                    # CLINK chain (9 layers)
rebis.materials.design_metamaterial(...)    # Materials science
rebis.imas.arrange_compound(...)            # IMASM compound signatures
rebis.gene.GeneticEngine(...)               # Gene imscriber
rebis.pipeline.imscribe_system(...)         # Auto-imscriber
rebis.shared.ORDINALS                       # Shared primitives & weights
rebis.demo.b4_lattice()                     # Quick demos
\end{verbatim}

\textbf{17 submodules} --- all import cleanly with zero stdout noise.

\begin{center}\rule{0.5\linewidth}{0.5pt}\end{center}

\begin{center}\rule{0.5\linewidth}{0.5pt}\end{center}

\hypertarget{overview}{%
\subsection{Overview}\label{overview}}

Red-Hot Rebis is the unified Imscribing Grammar bio/chem design
platform. It derives structural types for organisms, molecules,
proteins, and materials from the 12 primitives of the IG, then generates
physically actionable output files --- genome sequences, protein
structures, metabolic models, synthesis protocols --- directly from
those structural types.

The platform is organized around \textbf{seven core systems}:

\textbf{CLINK} --- 9-layer structural chain from subatomic quarks to
whole organism, each layer carrying an IG tuple. Every transition is
Frobenius-closed (\(\mu\circ\delta=\text{id}\)).

\textbf{IMASM} --- 12-token algebra (VINIT, TANCH, AFWD, AREV, CLINK,
IMSCRIB, FSPLIT, FFUSE, EVALT, EVALF, ENGAGR, IFIX) whose arrangement
classes map onto IG structural types. Now extended with the
\textbf{IMASM Compound Pipeline} for SMILES→IMASM encoding.

\textbf{SerpentRod} --- Frobenius morphism RNA→\{sequence+fold\} that
collapses gene → mature protein in one structural step.

\textbf{Paraconsistent Kernel (rhr\_p4rky)} --- Belnap FOUR logic,
paraconsistent abstract state machine, 64-codon B₄ lattice,
gene-to-protein pipeline, hadron/quark Belnap classification. 32 Python
modules (+ papers/).

\textbf{Gene Imscriber} --- IG-native genetic compiler --- structural
types to codon optimization, CRISPR guide design, chimera design,
Frobenius-verified editing.

\textbf{IMASM Compound Pipeline} --- SMILES→IMASM 8-token arrangement →
IG 12-tuple. 54 compounds registered, 190 FG patterns, cross-domain
analogy search across 4,027+ catalog entries.

\textbf{Crystal-Guided Molecular Discovery} --- navigates the
17,280,000-type Crystal of Types to reverse-engineer molecules from
target IG tuples. Discovered 5-nitro-bufotenin (DMT-⊙), the first
⊙-critical autocatalytic molecule.

\hypertarget{static-data-index.md}{%
\subsubsection{Static Data → INDEX.md}\label{static-data-index.md}}

Layer tables, IMASM canonical catalogs, materials listings, compound
catalog, and other reference data now lives in \textbf{INDEX.md} --- a
plain-text browsable reference. View it with:

\begin{verbatim}
less /home/mrnob0dy666/imsgct/red-hot_rebis/INDEX.md
\end{verbatim}

\hypertarget{the-12-primitives}{%
\subsubsection{The 12 Primitives}\label{the-12-primitives}}

Every structural entity carries a 12-primitive tuple. Each position is
one primitive; its value is a Shavian character from a 49-element set:

\begin{longtable}[]{@{}llll@{}}
\toprule
\begin{minipage}[b]{0.07\columnwidth}\raggedright
\#\strut
\end{minipage} & \begin{minipage}[b]{0.24\columnwidth}\raggedright
Primitive\strut
\end{minipage} & \begin{minipage}[b]{0.13\columnwidth}\raggedright
Name\strut
\end{minipage} & \begin{minipage}[b]{0.44\columnwidth}\raggedright
Values (low→high)\strut
\end{minipage}\tabularnewline
\midrule
\endhead
\begin{minipage}[t]{0.07\columnwidth}\raggedright
1\strut
\end{minipage} & \begin{minipage}[t]{0.24\columnwidth}\raggedright
\(Ð\)\strut
\end{minipage} & \begin{minipage}[t]{0.13\columnwidth}\raggedright
Dimensionality\strut
\end{minipage} & \begin{minipage}[t]{0.44\columnwidth}\raggedright
𐑛 (0d) → 𐑨 (2d) → 𐑼 (∞-dim) → 𐑦 (imscriptive)\strut
\end{minipage}\tabularnewline
\begin{minipage}[t]{0.07\columnwidth}\raggedright
2\strut
\end{minipage} & \begin{minipage}[t]{0.24\columnwidth}\raggedright
\(Þ\)\strut
\end{minipage} & \begin{minipage}[t]{0.13\columnwidth}\raggedright
Topology\strut
\end{minipage} & \begin{minipage}[t]{0.44\columnwidth}\raggedright
𐑡 (network) → 𐑰 (inclusion) → 𐑥 (crossing) → 𐑶 (box product) → 𐑸
(imscriptive closure)\strut
\end{minipage}\tabularnewline
\begin{minipage}[t]{0.07\columnwidth}\raggedright
3\strut
\end{minipage} & \begin{minipage}[t]{0.24\columnwidth}\raggedright
\(Ř\)\strut
\end{minipage} & \begin{minipage}[t]{0.13\columnwidth}\raggedright
Coupling\strut
\end{minipage} & \begin{minipage}[t]{0.44\columnwidth}\raggedright
𐑩 (supervenience) → 𐑑 (categorical) → 𐑽 (adjoint) → 𐑾
(bidirectional)\strut
\end{minipage}\tabularnewline
\begin{minipage}[t]{0.07\columnwidth}\raggedright
4\strut
\end{minipage} & \begin{minipage}[t]{0.24\columnwidth}\raggedright
\(Φ\)\strut
\end{minipage} & \begin{minipage}[t]{0.13\columnwidth}\raggedright
Parity (symmetry)\strut
\end{minipage} & \begin{minipage}[t]{0.44\columnwidth}\raggedright
𐑗 (none) → 𐑿 (quantum) → 𐑬 (partial) → 𐑯 (full) → 𐑹
(Frobenius-special)\strut
\end{minipage}\tabularnewline
\begin{minipage}[t]{0.07\columnwidth}\raggedright
5\strut
\end{minipage} & \begin{minipage}[t]{0.24\columnwidth}\raggedright
\(ƒ\)\strut
\end{minipage} & \begin{minipage}[t]{0.13\columnwidth}\raggedright
Fidelity\strut
\end{minipage} & \begin{minipage}[t]{0.44\columnwidth}\raggedright
𐑱 (classical) → 𐑞 (thermal) → 𐑐 (quantum coherence)\strut
\end{minipage}\tabularnewline
\begin{minipage}[t]{0.07\columnwidth}\raggedright
6\strut
\end{minipage} & \begin{minipage}[t]{0.24\columnwidth}\raggedright
\(Ç\)\strut
\end{minipage} & \begin{minipage}[t]{0.13\columnwidth}\raggedright
Kinetics\strut
\end{minipage} & \begin{minipage}[t]{0.44\columnwidth}\raggedright
𐑘 (MBL-frozen) → 𐑤 (trapped) → 𐑧 (near-equilibrium) → 𐑪 (moderate) → 𐑺
(fast/driven)\strut
\end{minipage}\tabularnewline
\begin{minipage}[t]{0.07\columnwidth}\raggedright
7\strut
\end{minipage} & \begin{minipage}[t]{0.24\columnwidth}\raggedright
\(Γ\)\strut
\end{minipage} & \begin{minipage}[t]{0.13\columnwidth}\raggedright
Cardinality\strut
\end{minipage} & \begin{minipage}[t]{0.44\columnwidth}\raggedright
𐑚 (local) → 𐑔 (mesoscale) → 𐑲 (maximal)\strut
\end{minipage}\tabularnewline
\begin{minipage}[t]{0.07\columnwidth}\raggedright
8\strut
\end{minipage} & \begin{minipage}[t]{0.24\columnwidth}\raggedright
\(ɢ\)\strut
\end{minipage} & \begin{minipage}[t]{0.13\columnwidth}\raggedright
Composition\strut
\end{minipage} & \begin{minipage}[t]{0.44\columnwidth}\raggedright
𐑝 (conjunctive) → 𐑜 (disjunctive) → 𐑠 (sequential) → 𐑵 (broadcast)\strut
\end{minipage}\tabularnewline
\begin{minipage}[t]{0.07\columnwidth}\raggedright
9\strut
\end{minipage} & \begin{minipage}[t]{0.24\columnwidth}\raggedright
\(⊙\)\strut
\end{minipage} & \begin{minipage}[t]{0.13\columnwidth}\raggedright
Criticality\strut
\end{minipage} & \begin{minipage}[t]{0.44\columnwidth}\raggedright
𐑢 (sub-critical) → ⊙ (critical / self-modeling) → 𐑮 (complex-plane) → 𐑻
(exceptional point) → 𐑣 (supercritical)\strut
\end{minipage}\tabularnewline
\begin{minipage}[t]{0.07\columnwidth}\raggedright
10\strut
\end{minipage} & \begin{minipage}[t]{0.24\columnwidth}\raggedright
\(Ħ\)\strut
\end{minipage} & \begin{minipage}[t]{0.13\columnwidth}\raggedright
Chirality\strut
\end{minipage} & \begin{minipage}[t]{0.44\columnwidth}\raggedright
𐑓 (memoryless) → 𐑒 (1-step) → 𐑖 (2-step) → 𐑫 (eternal)\strut
\end{minipage}\tabularnewline
\begin{minipage}[t]{0.07\columnwidth}\raggedright
11\strut
\end{minipage} & \begin{minipage}[t]{0.24\columnwidth}\raggedright
\(Σ\)\strut
\end{minipage} & \begin{minipage}[t]{0.13\columnwidth}\raggedright
Stoichiometry\strut
\end{minipage} & \begin{minipage}[t]{0.44\columnwidth}\raggedright
𐑙 (1:1) → 𐑕 (many identical) → 𐑳 (many heterogeneous)\strut
\end{minipage}\tabularnewline
\begin{minipage}[t]{0.07\columnwidth}\raggedright
12\strut
\end{minipage} & \begin{minipage}[t]{0.24\columnwidth}\raggedright
\(Ω\)\strut
\end{minipage} & \begin{minipage}[t]{0.13\columnwidth}\raggedright
Winding\strut
\end{minipage} & \begin{minipage}[t]{0.44\columnwidth}\raggedright
𐑷 (trivial) → 𐑴 (ℤ₂ parity) → 𐑭 (ℤ integer) → 𐑟 (non-Abelian)\strut
\end{minipage}\tabularnewline
\bottomrule
\end{longtable}

The standard tuple format:
\[\langle Ð;\ Þ;\ Ř;\ Φ;\ ƒ;\ Ç;\ Γ;\ ɢ;\ ⊙;\ Ħ;\ Σ;\ Ω \rangle\]

\begin{center}\rule{0.5\linewidth}{0.5pt}\end{center}

\hypertarget{installation}{%
\subsection{Installation}\label{installation}}

\begin{verbatim}
cd /home/mrnob0dy666/imsgct/red-hot_rebis

# Install the package and its dependencies (always use uv, never pip)
uv pip install -e .

# Install bridge packages used by scripts
uv pip install -e /home/mrnob0dy666/imsgct/omonad_OS
uv pip install -e /home/mrnob0dy666/imsgct/imasmic_core

# Optional: RDKit for molecular design (crystal designer, odot_finder)
uv pip install rdkit-pypi
\end{verbatim}

Verify everything is wired up:

\begin{verbatim}
python3 rebis.py verify
\end{verbatim}

Expected: 15+ lines all showing \texttt{✅}.

\begin{center}\rule{0.5\linewidth}{0.5pt}\end{center}

\hypertarget{quick-start}{%
\subsection{Quick Start}\label{quick-start}}

\begin{verbatim}
# System health — all modules with size and status
python3 rebis.py status

# Full verification (all imports, Frobenius closure)
python3 rebis.py verify

# List all 39 discoverable runnable targets
python3 rebis.py run list

# Run SerpentRod v5 on built-in test cases
python3 rebis.py run serpentrod

# Run the paraconsistent genetics test suite
python3 rebis.py run test_genetics

# CH3MPILER retrosynthetic compiler
python3 rebis.py run ch3mpiler --smiles "CC(=O)OC1=CC=CC=C1C(=O)O"  # Exact SMILES input
python3 rebis.py run ch3mpiler --target aspirin                                   # Name resolve
python3 rebis.py run ch3mpiler --target aspirin --smiles "CC(=O)OC1=CC=CC=C1C(=O)O"  # Both
python3 rebis.py run ch3mpiler --help

# Ch3mpiler + SerpentRod catalytic site design (CAS lookup)
python3 rebis.py run ch3mpiler_serpentrod_pipeline --cas 58-08-2

# IMASM Compound Pipeline — encode a molecule
python3 rebis.py imas compound --smiles "CC(=O)OC1=CC=CC=C1C(=O)O"

# IMASM Compound Pipeline — find cross-domain analogies
python3 rebis.py imas analogies --smiles "CN(C)CCC1=CNC2=CC=CC=C12" --limit 10

# IMASM Compound Pipeline — register a compound
python3 rebis.py imas register --smiles "CN1C=NC2=C1C(=O)N(C(=O)N2C)C" --name caffeine

# IMASM Reaction Analysis
python3 rebis.py imas reaction --smiles "CC(=O)C1=CC=CC=C1" --product "CC(C1=CC=CC=C1)(C)O"

# Comprehensive protein & genetics stress test (34 tests)
python3 rebis.py run stress_test_proteins

# Materials stress test (26 tests)
python3 materials/stress_test_materials.py

# Generate a full mammal organism design package
python3 rebis.py pipeline actionable

# Browse static reference data (CLINK layers, IMASM canonicals, compound catalog)
less INDEX.md
\end{verbatim}

\begin{center}\rule{0.5\linewidth}{0.5pt}\end{center}

\hypertarget{command-reference}{%
\subsection{Command Reference}\label{command-reference}}

\hypertarget{status}{%
\subsubsection{\texorpdfstring{\texttt{status}}{status}}\label{status}}

Shows all platform modules with file size and health check.

\begin{verbatim}
python3 rebis.py status
\end{verbatim}

Output includes: SerpentRod, CH3MPILER, Pipeline, Gene Imscriber, CLINK,
IMASM modules, \textbf{Compound Pipeline}, Materials, Therapeutics,
Paraconsistent Kernel (rhr\_p4rky), Popular Protein --- each with file
size and ✅/❌.

\begin{center}\rule{0.5\linewidth}{0.5pt}\end{center}

\hypertarget{verify}{%
\subsubsection{\texorpdfstring{\texttt{verify}}{verify}}\label{verify}}

Imports every module and reports pass/fail. Includes Frobenius closure
checks where applicable. Run this after any install or code change.

\begin{verbatim}
python3 rebis.py verify
\end{verbatim}

\begin{center}\rule{0.5\linewidth}{0.5pt}\end{center}

\hypertarget{run}{%
\subsubsection{\texorpdfstring{\texttt{run}}{run}}\label{run}}

Run one of 39 discoverable scripts or modules. Use \texttt{run\ list} to
see all targets.

\begin{verbatim}
python3 rebis.py run <target> [args...]
python3 rebis.py run list           # Show all 39 discoverable targets
\end{verbatim}

\hypertarget{core-module-targets}{%
\paragraph{Core module targets}\label{core-module-targets}}

\begin{longtable}[]{@{}lll@{}}
\toprule
\begin{minipage}[b]{0.24\columnwidth}\raggedright
Target\strut
\end{minipage} & \begin{minipage}[b]{0.24\columnwidth}\raggedright
Module\strut
\end{minipage} & \begin{minipage}[b]{0.43\columnwidth}\raggedright
What it does\strut
\end{minipage}\tabularnewline
\midrule
\endhead
\begin{minipage}[t]{0.24\columnwidth}\raggedright
\texttt{serpentrod}\strut
\end{minipage} & \begin{minipage}[t]{0.24\columnwidth}\raggedright
\texttt{serpentrod/protein\_v5.py}\strut
\end{minipage} & \begin{minipage}[t]{0.43\columnwidth}\raggedright
SerpentRod v5 --- full signal peptide detection, cleavage, fragment
naming, PTMs, primitive spectrum. Runs built-in test cases (Human
Insulin, Proglucagon, Preproinsulin, Proopiomelanocortin, GLP-1
analog)\strut
\end{minipage}\tabularnewline
\begin{minipage}[t]{0.24\columnwidth}\raggedright
\texttt{serpentrod\_v2}\strut
\end{minipage} & \begin{minipage}[t]{0.24\columnwidth}\raggedright
\texttt{rhr\_p4rky/serpent\_rod\_v2.py}\strut
\end{minipage} & \begin{minipage}[t]{0.43\columnwidth}\raggedright
SerpentRod v2 --- 3D backbone from φ/ψ angles, geometry-based contacts,
energy scoring\strut
\end{minipage}\tabularnewline
\begin{minipage}[t]{0.24\columnwidth}\raggedright
\texttt{ch3mpiler}\strut
\end{minipage} & \begin{minipage}[t]{0.24\columnwidth}\raggedright
\texttt{ch3mpiler/compiler.py}\strut
\end{minipage} & \begin{minipage}[t]{0.43\columnwidth}\raggedright
Retrosynthetic disconnection compiler --- \texttt{-\/-smiles} for direct
SMILES, \texttt{-\/-target} for name resolution,
\texttt{-\/-retrosynthesis} for ranked disconnections with \textbf{exact
fragment SMILES}\strut
\end{minipage}\tabularnewline
\begin{minipage}[t]{0.24\columnwidth}\raggedright
\texttt{ch3mpiler\_serpentrod\_pipeline}\strut
\end{minipage} & \begin{minipage}[t]{0.24\columnwidth}\raggedright
\texttt{rhr\_p4rky/ch3mpiler\_serpentrod\_pipeline.py}\strut
\end{minipage} & \begin{minipage}[t]{0.43\columnwidth}\raggedright
Ch3mpiler + SerpentRod integrated pipeline --- catalytic site design
from SMILES or CAS\strut
\end{minipage}\tabularnewline
\begin{minipage}[t]{0.24\columnwidth}\raggedright
\texttt{gene}\strut
\end{minipage} & \begin{minipage}[t]{0.24\columnwidth}\raggedright
\texttt{gene\_imscriber/engine.py}\strut
\end{minipage} & \begin{minipage}[t]{0.43\columnwidth}\raggedright
Gene imscriber (structural type → codon optimization)\strut
\end{minipage}\tabularnewline
\begin{minipage}[t]{0.24\columnwidth}\raggedright
\texttt{test\_genetics}\strut
\end{minipage} & \begin{minipage}[t]{0.24\columnwidth}\raggedright
\texttt{test\_genetics.py}\strut
\end{minipage} & \begin{minipage}[t]{0.43\columnwidth}\raggedright
64-codon B4 lattice test suite\strut
\end{minipage}\tabularnewline
\begin{minipage}[t]{0.24\columnwidth}\raggedright
\texttt{stress\_test\_proteins}\strut
\end{minipage} & \begin{minipage}[t]{0.24\columnwidth}\raggedright
\texttt{stress\_test\_proteins.py}\strut
\end{minipage} & \begin{minipage}[t]{0.43\columnwidth}\raggedright
Comprehensive 34-test protein \& genetics stress test\strut
\end{minipage}\tabularnewline
\begin{minipage}[t]{0.24\columnwidth}\raggedright
\texttt{compare\_exact}\strut
\end{minipage} & \begin{minipage}[t]{0.24\columnwidth}\raggedright
\texttt{popular\_protein/compare\_exact\_phi\_psi.py}\strut
\end{minipage} & \begin{minipage}[t]{0.43\columnwidth}\raggedright
Exact φ/ψ angle comparison\strut
\end{minipage}\tabularnewline
\begin{minipage}[t]{0.24\columnwidth}\raggedright
\texttt{compare\_structures}\strut
\end{minipage} & \begin{minipage}[t]{0.24\columnwidth}\raggedright
\texttt{popular\_protein/compare\_structures.py}\strut
\end{minipage} & \begin{minipage}[t]{0.43\columnwidth}\raggedright
Structure-level comparison\strut
\end{minipage}\tabularnewline
\begin{minipage}[t]{0.24\columnwidth}\raggedright
\texttt{comprehensive\_comparison}\strut
\end{minipage} & \begin{minipage}[t]{0.24\columnwidth}\raggedright
\texttt{popular\_protein/comprehensive\_comparison.py}\strut
\end{minipage} & \begin{minipage}[t]{0.43\columnwidth}\raggedright
Full multi-metric protein report\strut
\end{minipage}\tabularnewline
\begin{minipage}[t]{0.24\columnwidth}\raggedright
\texttt{run\_gene\_pipeline}\strut
\end{minipage} & \begin{minipage}[t]{0.24\columnwidth}\raggedright
\texttt{rhr\_p4rky/run\_gene\_pipeline.py}\strut
\end{minipage} & \begin{minipage}[t]{0.43\columnwidth}\raggedright
Gene→protein 7-stage B4 pipeline\strut
\end{minipage}\tabularnewline
\begin{minipage}[t]{0.24\columnwidth}\raggedright
\texttt{antibody}\strut
\end{minipage} & \begin{minipage}[t]{0.24\columnwidth}\raggedright
\texttt{run\_antibody.py}\strut
\end{minipage} & \begin{minipage}[t]{0.43\columnwidth}\raggedright
Antibody CDR designer\strut
\end{minipage}\tabularnewline
\begin{minipage}[t]{0.24\columnwidth}\raggedright
\texttt{genetic\_code}\strut
\end{minipage} & \begin{minipage}[t]{0.24\columnwidth}\raggedright
\texttt{rhr\_p4rky/genetic\_code.py}\strut
\end{minipage} & \begin{minipage}[t]{0.43\columnwidth}\raggedright
64-codon Frobenius-verified table\strut
\end{minipage}\tabularnewline
\begin{minipage}[t]{0.24\columnwidth}\raggedright
\texttt{pschedelic}\strut
\end{minipage} & \begin{minipage}[t]{0.24\columnwidth}\raggedright
\texttt{psychedelic\_bridge.py}\strut
\end{minipage} & \begin{minipage}[t]{0.43\columnwidth}\raggedright
Compound intrinsics + coupling report\strut
\end{minipage}\tabularnewline
\begin{minipage}[t]{0.24\columnwidth}\raggedright
\texttt{iupac}\strut
\end{minipage} & \begin{minipage}[t]{0.24\columnwidth}\raggedright
\texttt{\_gen\_univ\_map.py}\strut
\end{minipage} & \begin{minipage}[t]{0.43\columnwidth}\raggedright
Diaschizic IUPAC generator\strut
\end{minipage}\tabularnewline
\begin{minipage}[t]{0.24\columnwidth}\raggedright
\texttt{mito}\strut
\end{minipage} & \begin{minipage}[t]{0.24\columnwidth}\raggedright
\texttt{mito\_pipeline.py}\strut
\end{minipage} & \begin{minipage}[t]{0.43\columnwidth}\raggedright
Mitochondrial gene pipeline\strut
\end{minipage}\tabularnewline
\begin{minipage}[t]{0.24\columnwidth}\raggedright
\texttt{mito\_v2}\strut
\end{minipage} & \begin{minipage}[t]{0.24\columnwidth}\raggedright
\texttt{mito\_pipeline\_v2.py}\strut
\end{minipage} & \begin{minipage}[t]{0.43\columnwidth}\raggedright
Mitochondrial gene pipeline v2\strut
\end{minipage}\tabularnewline
\begin{minipage}[t]{0.24\columnwidth}\raggedright
\texttt{mito\_v3}\strut
\end{minipage} & \begin{minipage}[t]{0.24\columnwidth}\raggedright
\texttt{mito\_pipeline\_v3.py}\strut
\end{minipage} & \begin{minipage}[t]{0.43\columnwidth}\raggedright
Mitochondrial gene pipeline v3\strut
\end{minipage}\tabularnewline
\begin{minipage}[t]{0.24\columnwidth}\raggedright
\texttt{antibody\_designer}\strut
\end{minipage} & \begin{minipage}[t]{0.24\columnwidth}\raggedright
\texttt{rhr\_p4rky/antibody\_designer.py}\strut
\end{minipage} & \begin{minipage}[t]{0.43\columnwidth}\raggedright
Advanced antibody designer\strut
\end{minipage}\tabularnewline
\begin{minipage}[t]{0.24\columnwidth}\raggedright
\texttt{pdb\_validator}\strut
\end{minipage} & \begin{minipage}[t]{0.24\columnwidth}\raggedright
\texttt{rhr\_p4rky/pdb\_validator.py}\strut
\end{minipage} & \begin{minipage}[t]{0.43\columnwidth}\raggedright
PDB structure validation\strut
\end{minipage}\tabularnewline
\begin{minipage}[t]{0.24\columnwidth}\raggedright
\texttt{hadron\_belnap}\strut
\end{minipage} & \begin{minipage}[t]{0.24\columnwidth}\raggedright
\texttt{rhr\_p4rky/hadron\_belnap.py}\strut
\end{minipage} & \begin{minipage}[t]{0.43\columnwidth}\raggedright
Hadronic Belnap classification\strut
\end{minipage}\tabularnewline
\begin{minipage}[t]{0.24\columnwidth}\raggedright
\texttt{exotic\_hadron\_belnap}\strut
\end{minipage} & \begin{minipage}[t]{0.24\columnwidth}\raggedright
\texttt{rhr\_p4rky/exotic\_hadron\_belnap.py}\strut
\end{minipage} & \begin{minipage}[t]{0.43\columnwidth}\raggedright
Exotic hadron classification\strut
\end{minipage}\tabularnewline
\begin{minipage}[t]{0.24\columnwidth}\raggedright
\texttt{quark\_belnap}\strut
\end{minipage} & \begin{minipage}[t]{0.24\columnwidth}\raggedright
\texttt{rhr\_p4rky/quark\_belnap.py}\strut
\end{minipage} & \begin{minipage}[t]{0.43\columnwidth}\raggedright
Quark Belnap analysis\strut
\end{minipage}\tabularnewline
\bottomrule
\end{longtable}

\hypertarget{clink}{%
\subsubsection{\texorpdfstring{\texttt{clink}}{clink}}\label{clink}}

Inspect the CLINK chain --- 9 layers from frustrated quarks to whole
organisms.

\begin{verbatim}
python3 rebis.py clink <subcommand> [options]
\end{verbatim}

\hypertarget{clink-layer-name_or_index}{%
\paragraph{\texorpdfstring{\texttt{clink\ layer\ \textless{}name\_or\_index\textgreater{}}}{clink layer \textless name\_or\_index\textgreater{}}}\label{clink-layer-name_or_index}}

Show the full profile of one CLINK layer --- tuple, tier, Frobenius
status, bridges, promotion paths.

\begin{verbatim}
python3 rebis.py clink layer 8          # Organism
python3 rebis.py clink layer cell       # Fuzzy match → Cell layer
python3 rebis.py clink layer L3         # Molecule layer
\end{verbatim}

Output includes: index, name, tier, distance from previous, IG tuple
(Shavian glyphs), primitive-by-primitive description, and nearest tool
bridge.

\hypertarget{clink-bridge-component-target}{%
\paragraph{\texorpdfstring{\texttt{clink\ bridge\ \textless{}component\textgreater{}\ \textless{}target\textgreater{}}}{clink bridge \textless component\textgreater{} \textless target\textgreater{}}}\label{clink-bridge-component-target}}

Show the promotion path from a component tool to a CLINK target layer.

\begin{verbatim}
python3 rebis.py clink bridge serpentrod 8
python3 rebis.py clink bridge ch3mpiler 3
\end{verbatim}

Shows: distance, primitive-level promotion steps, Frobenius status at
each hop.

\textbf{Static reference:} \texttt{less\ INDEX.md} §1 (9-layer table,
tuples, tiers, Frobenius status, chain distances, ZFC\_fe distance,
component bridge attachments).

\begin{center}\rule{0.5\linewidth}{0.5pt}\end{center}

\hypertarget{pipeline}{%
\subsubsection{\texorpdfstring{\texttt{pipeline}}{pipeline}}\label{pipeline}}

Orchestrate the full CLINK pipeline --- design a whole organism from
structural types alone.

\begin{verbatim}
python3 rebis.py pipeline <subcommand> [options]
\end{verbatim}

\hypertarget{pipeline-bridges}{%
\paragraph{\texorpdfstring{\texttt{pipeline\ bridges}}{pipeline bridges}}\label{pipeline-bridges}}

Show which tool bridges are available (✅) or missing (❌).

\begin{verbatim}
python3 rebis.py pipeline bridges
\end{verbatim}

Current bridges: serpentrod ✅, ch3mpiler ✅, gene\_imscriber ✅,
materials\_sim ✅, frobenius\_chemo ✅, ouroboric\_pill ✅,
critical\_metamaterial ✅, neurotrophic ✅, biology\_sim ❌,
ouroboric\_telomere ❌.

\hypertarget{pipeline-ground-up}{%
\paragraph{\texorpdfstring{\texttt{pipeline\ ground-up}}{pipeline ground-up}}\label{pipeline-ground-up}}

Full L0→L8 synthesis run. Traverses all 9 layers, applies layer
designers, reports status per transition.

\begin{verbatim}
python3 rebis.py pipeline ground-up
\end{verbatim}

Output: per-layer transition table with distance, primitive promotions,
tool used, and ✅/❌. Total distance and total promotions at end.

\hypertarget{pipeline-from-layer-start-end}{%
\paragraph{\texorpdfstring{\texttt{pipeline\ from-layer\ \textless{}start\textgreater{}\ \textless{}end\textgreater{}}}{pipeline from-layer \textless start\textgreater{} \textless end\textgreater{}}}\label{pipeline-from-layer-start-end}}

Partial run starting from a specific layer.

\begin{verbatim}
python3 rebis.py pipeline from-layer 3 8   # molecule → organism
python3 rebis.py pipeline from-layer 4 8   # cell → organism
\end{verbatim}

Useful when you have an existing molecular design and want to take it up
to organism level.

\hypertarget{pipeline-actionable---organism-type}{%
\paragraph{\texorpdfstring{\texttt{pipeline\ actionable\ {[}-\/-organism\ \textless{}type\textgreater{}{]}}}{pipeline actionable {[}-\/-organism \textless type\textgreater{]}}}\label{pipeline-actionable---organism-type}}

Generate a complete, physically actionable organism design package.

\begin{verbatim}
python3 rebis.py pipeline actionable
python3 rebis.py pipeline actionable --organism human
\end{verbatim}

Produces a directory of files in
\texttt{clink/datasets/organism\_designs/organism\_\textless{}type\textgreater{}\_actionable/}:

\begin{longtable}[]{@{}ll@{}}
\toprule
\begin{minipage}[b]{0.40\columnwidth}\raggedright
File\strut
\end{minipage} & \begin{minipage}[b]{0.54\columnwidth}\raggedright
Action\strut
\end{minipage}\tabularnewline
\midrule
\endhead
\begin{minipage}[t]{0.40\columnwidth}\raggedright
\texttt{genome.fasta}\strut
\end{minipage} & \begin{minipage}[t]{0.54\columnwidth}\raggedright
Order DNA synthesis from Twist/IDT/GenScript\strut
\end{minipage}\tabularnewline
\begin{minipage}[t]{0.40\columnwidth}\raggedright
\texttt{genome.gb}\strut
\end{minipage} & \begin{minipage}[t]{0.54\columnwidth}\raggedright
Load into Benchling/SnapGene\strut
\end{minipage}\tabularnewline
\begin{minipage}[t]{0.40\columnwidth}\raggedright
\texttt{construct.sbol}\strut
\end{minipage} & \begin{minipage}[t]{0.54\columnwidth}\raggedright
Exchange with synthetic biology repositories\strut
\end{minipage}\tabularnewline
\begin{minipage}[t]{0.40\columnwidth}\raggedright
\texttt{codon\_usage.csv}\strut
\end{minipage} & \begin{minipage}[t]{0.54\columnwidth}\raggedright
B4-derived codon table\strut
\end{minipage}\tabularnewline
\begin{minipage}[t]{0.40\columnwidth}\raggedright
\texttt{protein.fasta}\strut
\end{minipage} & \begin{minipage}[t]{0.54\columnwidth}\raggedright
Order peptide synthesis\strut
\end{minipage}\tabularnewline
\begin{minipage}[t]{0.40\columnwidth}\raggedright
\texttt{protein\_coords.pdb}\strut
\end{minipage} & \begin{minipage}[t]{0.54\columnwidth}\raggedright
View in PyMOL/ChimeraX (PDB v3.3-compliant: SEQRES, TER, correct atom
names)\strut
\end{minipage}\tabularnewline
\begin{minipage}[t]{0.40\columnwidth}\raggedright
\texttt{serpentrod\_classification.json}\strut
\end{minipage} & \begin{minipage}[t]{0.54\columnwidth}\raggedright
Primitive activations per fragment\strut
\end{minipage}\tabularnewline
\begin{minipage}[t]{0.40\columnwidth}\raggedright
\texttt{molecules.smi}\strut
\end{minipage} & \begin{minipage}[t]{0.54\columnwidth}\raggedright
SMILES --- order from vendor\strut
\end{minipage}\tabularnewline
\begin{minipage}[t]{0.40\columnwidth}\raggedright
\texttt{retro\_pathways.json}\strut
\end{minipage} & \begin{minipage}[t]{0.54\columnwidth}\raggedright
CH3MPILER retrosynthesis\strut
\end{minipage}\tabularnewline
\begin{minipage}[t]{0.40\columnwidth}\raggedright
\texttt{metabolic\_model.xml}\strut
\end{minipage} & \begin{minipage}[t]{0.54\columnwidth}\raggedright
Load into COBRApy for FBA\strut
\end{minipage}\tabularnewline
\begin{minipage}[t]{0.40\columnwidth}\raggedright
\texttt{plasmid.gb}\strut
\end{minipage} & \begin{minipage}[t]{0.54\columnwidth}\raggedright
GenBank construct\strut
\end{minipage}\tabularnewline
\begin{minipage}[t]{0.40\columnwidth}\raggedright
\texttt{growth\_media.txt}\strut
\end{minipage} & \begin{minipage}[t]{0.54\columnwidth}\raggedright
Prepare media per formulation\strut
\end{minipage}\tabularnewline
\begin{minipage}[t]{0.40\columnwidth}\raggedright
\texttt{organoid\_protocol.md}\strut
\end{minipage} & \begin{minipage}[t]{0.54\columnwidth}\raggedright
Lab-ready protocol\strut
\end{minipage}\tabularnewline
\begin{minipage}[t]{0.40\columnwidth}\raggedright
\texttt{mitosis\_assay\_protocol.md}\strut
\end{minipage} & \begin{minipage}[t]{0.54\columnwidth}\raggedright
Lab-ready protocol\strut
\end{minipage}\tabularnewline
\begin{minipage}[t]{0.40\columnwidth}\raggedright
\texttt{physiological\_params.csv}\strut
\end{minipage} & \begin{minipage}[t]{0.54\columnwidth}\raggedright
Homeostatic setpoints\strut
\end{minipage}\tabularnewline
\begin{minipage}[t]{0.40\columnwidth}\raggedright
\texttt{whole\_genome\_spec.json}\strut
\end{minipage} & \begin{minipage}[t]{0.54\columnwidth}\raggedright
Complete genome specification\strut
\end{minipage}\tabularnewline
\bottomrule
\end{longtable}

Current designed organisms: mammal, human, human\_gills,
human\_photosynthetic, treople.

\begin{center}\rule{0.5\linewidth}{0.5pt}\end{center}

\hypertarget{materials}{%
\subsubsection{\texorpdfstring{\texttt{materials}}{materials}}\label{materials}}

Design novel IG-typed materials. All materials are derived from their
12-primitive structural types. Five dynamic subcommands; static
reference data (material catalog, Sophick Mercury, Eagle Cycle, gap
primitives) is in \textbf{INDEX.md}.

\begin{verbatim}
python3 rebis.py materials <subcommand> [options]
\end{verbatim}

\hypertarget{materials-forge---name-material---all}{%
\paragraph{\texorpdfstring{\texttt{materials\ forge\ -\/-name\ \textless{}material\textgreater{}\ {[}-\/-all{]}}}{materials forge -\/-name \textless material\textgreater{} {[}-\/-all{]}}}\label{materials-forge---name-material---all}}

Generate the complete material design file.

\begin{verbatim}
python3 rebis.py materials forge --name frobenius_composite
python3 rebis.py materials forge --name I_Dialetheic_Bootstrap   # From IMASM canonical
python3 rebis.py materials forge --all     # Forge all 8 predefined materials
\end{verbatim}

\hypertarget{materials-frobenius}{%
\paragraph{\texorpdfstring{\texttt{materials\ frobenius}}{materials frobenius}}\label{materials-frobenius}}

Simulate Frobenius closure verification of the composite --- cyclic
load/heal protocol showing \(\|\mu\delta-\text{id}\|\) per cycle.

\begin{verbatim}
python3 rebis.py materials frobenius
\end{verbatim}

\hypertarget{materials-ouroboric}{%
\paragraph{\texorpdfstring{\texttt{materials\ ouroboric}}{materials ouroboric}}\label{materials-ouroboric}}

Simulate ouroboric crack healing --- crack propagation + self-healing
cycles, reports fatigue life and topological invariant preservation.

\begin{verbatim}
python3 rebis.py materials ouroboric
\end{verbatim}

\hypertarget{materials-sophick---name-target}{%
\paragraph{\texorpdfstring{\texttt{materials\ sophick\ -\/-name\ \textless{}target\textgreater{}}}{materials sophick -\/-name \textless target\textgreater{}}}\label{materials-sophick---name-target}}

Eagle Cycle Protocol --- prepares an \(\text{O}_{\infty}\) Sophick
Mercury substrate from O₂ ouroboric materials. Requires
\texttt{-\/-name}.

\begin{verbatim}
python3 rebis.py materials sophick --name eagle_9_sophick
python3 rebis.py materials sophick --name cliff               # Frobenius Cliff analysis
python3 rebis.py materials sophick --name bridge              # IMASM→Eagle bridge report
\end{verbatim}

\hypertarget{materials-exactor---name-target}{%
\paragraph{\texorpdfstring{\texttt{materials\ exactor\ -\/-name\ \textless{}target\textgreater{}}}{materials exactor -\/-name \textless target\textgreater{}}}\label{materials-exactor---name-target}}

Explains the thermodynamic ceiling: continuous Eagle preparation reaches
its limit, and exact Frobenius closure requires discrete topological
protection. Requires \texttt{-\/-name}.

\begin{verbatim}
python3 rebis.py materials exactor --name diagnose            # Category error diagnosis
python3 rebis.py materials exactor --name close               # Close Frobenius gap
python3 rebis.py materials exactor --name pathways            # List all exactor pathways
\end{verbatim}

\hypertarget{static-reference-less-index.md-3-8-predefined-materials-imasm-canonicals-as-materials-sophick-mercury-pathway-frobenius-exactor-gap-primitives.}{%
\subsection{\texorpdfstring{\textbf{Static reference:}
\texttt{less\ INDEX.md} §3 (8 predefined materials, IMASM canonicals as
materials, Sophick Mercury pathway, Frobenius Exactor gap
primitives).}{Static reference: less INDEX.md §3 (8 predefined materials, IMASM canonicals as materials, Sophick Mercury pathway, Frobenius Exactor gap primitives).}}\label{static-reference-less-index.md-3-8-predefined-materials-imasm-canonicals-as-materials-sophick-mercury-pathway-frobenius-exactor-gap-primitives.}}

\hypertarget{imas}{%
\subsubsection{\texorpdfstring{\texttt{imas}}{imas}}\label{imas}}

IMASM arrangement analysis --- the 12-token algebra over IG structural
types. Three existing subcommands (\texttt{bridge}, \texttt{hunt},
\texttt{energy}) plus \textbf{three new compound pipeline subcommands}
(\texttt{compound}, \texttt{analogies}, \texttt{reaction},
\texttt{register}).

\begin{verbatim}
python3 rebis.py imas <subcommand> [options]
\end{verbatim}

\hypertarget{imas-bridge---canonical-name}{%
\paragraph{\texorpdfstring{\texttt{imas\ bridge\ {[}-\/-canonical\ \textless{}name\textgreater{}{]}}}{imas bridge {[}-\/-canonical \textless name\textgreater{]}}}\label{imas-bridge---canonical-name}}

Detailed profile of one IMASM canonical --- full IG tuple
primitive-by-primitive, nearest CLINK layer, distance table to all 9
layers.

\begin{verbatim}
python3 rebis.py imas bridge --canonical I_Dialetheic_Bootstrap
python3 rebis.py imas bridge --canonical I_Dialetheic_Bootstrap,VII_Parakernel  # Multi
\end{verbatim}

Default: \texttt{I\_Dialetheic\_Bootstrap}.

\hypertarget{imas-hunt---samples-n}{%
\paragraph{\texorpdfstring{\texttt{imas\ hunt\ {[}-\/-samples\ \textless{}n\textgreater{}{]}}}{imas hunt {[}-\/-samples \textless n\textgreater{]}}}\label{imas-hunt---samples-n}}

Monte Carlo Frobenius density estimation over the 12-token sequence
space. Reports probability of each Frobenius class and generates
examples.

\begin{verbatim}
python3 rebis.py imas hunt --samples 100000
python3 rebis.py imas hunt --samples 1000000          # Larger sample
\end{verbatim}

Output: \(p_{\text{frobenius\_pair}} \approx 0.236\),
\(p_{\text{proper\_frobenius}} \approx 0.139\),
\(p_{\text{dialetheia\_complete}} \approx 0.105\).

\hypertarget{imas-energy---canonical-name---layer-idx}{%
\paragraph{\texorpdfstring{\texttt{imas\ energy\ {[}-\/-canonical\ \textless{}name\textgreater{}{]}\ {[}-\/-layer\ \textless{}idx\textgreater{}{]}}}{imas energy {[}-\/-canonical \textless name\textgreater{]} {[}-\/-layer \textless idx\textgreater{]}}}\label{imas-energy---canonical-name---layer-idx}}

Structural activation energy from an IMASM canonical to a CLINK target
layer.

\begin{verbatim}
python3 rebis.py imas energy --canonical I_Dialetheic_Bootstrap --layer L8_Organism
python3 rebis.py imas energy --canonical V_Linear_Chain --layer L0_FrustratedBelnap5
\end{verbatim}

Default: \texttt{I\_Dialetheic\_Bootstrap} → L8 (Whole Organism).

\hypertarget{imas-compound---smiles-smiles---name-name---json}{%
\paragraph{\texorpdfstring{\texttt{imas\ compound\ -\/-smiles\ \textless{}SMILES\textgreater{}\ {[}-\/-name\ \textless{}name\textgreater{}{]}\ {[}-\/-json{]}}}{imas compound -\/-smiles \textless SMILES\textgreater{} {[}-\/-name \textless name\textgreater{]} {[}-\/-json{]}}}\label{imas-compound---smiles-smiles---name-name---json}}

\textbf{NEW --- Encode a molecule as an IMASM arrangement.} The compound
pipeline detects functional groups (from 190 SMARTS patterns), builds an
8-token arrangement, computes the StructuralFingerprint, and resolves to
an IG 12-tuple.

\begin{verbatim}
# Aspirin — print arrangement + fingerprint + IG tuple
python3 rebis.py imas compound --smiles "CC(=O)OC1=CC=CC=C1C(=O)O"

# Caffeine — JSON output for programmatic use
python3 rebis.py imas compound --smiles "CN1C=NC2=C1C(=O)N(C(=O)N2C)C" --json

# DMT — with name annotation
python3 rebis.py imas compound --smiles "CN(C)CCC1=CNC2=CC=CC=C12" --name "DMT"
\end{verbatim}

Output includes: detected FGs, 8-token IMASM arrangement,
StructuralFingerprint fields (self\_ref, dialetheia\_complete,
frobenius\_order, period, diversity, sigmas), IG 12-tuple with
per-primitive description, tier assessment.

\hypertarget{imas-analogies---smiles-smiles---limit-n}{%
\paragraph{\texorpdfstring{\texttt{imas\ analogies\ -\/-smiles\ \textless{}SMILES\textgreater{}\ {[}-\/-limit\ \textless{}n\textgreater{}{]}}}{imas analogies -\/-smiles \textless SMILES\textgreater{} {[}-\/-limit \textless n\textgreater{]}}}\label{imas-analogies---smiles-smiles---limit-n}}

\textbf{NEW --- Find cross-domain structural analogies.} Given a SMILES,
compute its IG type and search the entire catalog for nearest neighbors
--- across materials, consciousness states, languages, mathematical
theorems, and more.

\begin{verbatim}
# DMT's 10 nearest neighbors
python3 rebis.py imas analogies --smiles "CN(C)CCC1=CNC2=CC=CC=C12" --limit 10

# Water's nearest neighbors
python3 rebis.py imas analogies --smiles "O" --limit 15
\end{verbatim}

Output: ranked list of catalog entries with structural distances.
Typical results include consciousness states (Kalachakra Tantra, d=4),
mathematical theorems (Fermat's Last Theorem, d=4), languages (Codex
Vienna, d=5), materials (supercritical water, d=5).

\hypertarget{imas-reaction---smiles-smiles---product-smiles}{%
\paragraph{\texorpdfstring{\texttt{imas\ reaction\ -\/-smiles\ \textless{}SMILES\textgreater{}\ -\/-product\ \textless{}SMILES\textgreater{}}}{imas reaction -\/-smiles \textless SMILES\textgreater{} -\/-product \textless SMILES\textgreater{}}}\label{imas-reaction---smiles-smiles---product-smiles}}

\textbf{NEW --- Analyze a reaction as an IMASM transition.} Computes
reactant and product IMASM arrangements, identifies the token delta
(what changed), and classifies the reaction.

\begin{verbatim}
# Grignard: acetophenone → 2-phenyl-2-propanol
python3 rebis.py imas reaction --smiles "CC(=O)C1=CC=CC=C1" --product "CC(C1=CC=CC=C1)(C)O"
\end{verbatim}

Output includes: reactant arrangement, product arrangement, IMASM delta
(AFWD→AREV), reaction class (Redox\_reduction), named reaction
confidence scores.

\hypertarget{imas-register---smiles-smiles---name-name}{%
\paragraph{\texorpdfstring{\texttt{imas\ register\ -\/-smiles\ \textless{}SMILES\textgreater{}\ -\/-name\ \textless{}name\textgreater{}}}{imas register -\/-smiles \textless SMILES\textgreater{} -\/-name \textless name\textgreater{}}}\label{imas-register---smiles-smiles---name-name}}

\textbf{NEW --- Register a compound in the IG catalog.} Encodes the
molecule, computes its IG tuple, and registers it in both the
COMPOUND\_DATABASE and the shared IG\_catalog.json for future lookups.

\begin{verbatim}
python3 rebis.py imas register --smiles "CN1C=NC2=C1C(=O)N(C(=O)N2C)C" --name "caffeine"
python3 rebis.py imas register --smiles "CC(=O)OC1=CC=CC=C1C(=O)O" --name "aspirin"
\end{verbatim}

\textbf{Static reference:} \texttt{less\ INDEX.md} §4 (compound catalog
with statistics, distinct IG types, cross-domain analogies table).

\begin{center}\rule{0.5\linewidth}{0.5pt}\end{center}

\hypertarget{scripts}{%
\subsubsection{\texorpdfstring{\texttt{scripts}}{scripts}}\label{scripts}}

Manage and run standalone scripts.

\begin{verbatim}
python3 rebis.py scripts list              # List all 14 scripts with line counts
python3 rebis.py scripts run <name>        # Run a script by name (without .py)
\end{verbatim}

\begin{verbatim}
python3 rebis.py scripts run mito_pipeline
python3 rebis.py scripts run omonad_bridge
python3 rebis.py scripts run run_pdb_validation
\end{verbatim}

Note: \texttt{run\ mito}, \texttt{run\ antibody},
\texttt{run\ psychedelic}, \texttt{run\ iupac} are convenience aliases
for \texttt{scripts\ run}.

\begin{center}\rule{0.5\linewidth}{0.5pt}\end{center}

\hypertarget{crystal-guided-molecular-discovery}{%
\subsection{Crystal-Guided Molecular
Discovery}\label{crystal-guided-molecular-discovery}}

The \textbf{Crystal-Guided Molecular Discovery Engine}
(\texttt{imas/molecular\_crystal\_designer.py}, 700+ lines) navigates
the Imscribing Grammar's Crystal of Types (17,280,000 structural types)
to discover molecules with target structural properties. It is the
\textbf{inverse} of the SMILES→IMASM→IG pipeline.

\hypertarget{how-it-works}{%
\subsubsection{How It Works}\label{how-it-works}}

\begin{verbatim}
Target IG tuple → [ig_to_fp_constraints] → FP constraints → [arrangement candidates]
  → [molecular_suggestion] → candidate SMILES → [validate with RDKit]
\end{verbatim}

The discovery protocol: 1. Encode molecule → IMASM arrangement → IG
12-tuple 2. Navigate crystal neighborhood → find target types 3. Target
⊙: requires self\_ref=True AND dialetheia\_complete=True 4. Design
constraints: EVALT + EVALF + ENGAGR in 8-token arrangement 5. Minimize
competing high-priority tokens (phenol \textgreater{} carboxylic acid)
6. Generate candidate SMILES → validate → register

\hypertarget{discovery-5-nitro-bufotenin-dmt-}{%
\subsubsection{Discovery: 5-Nitro-Bufotenin
(DMT-⊙)}\label{discovery-5-nitro-bufotenin-dmt-}}

The first discovery was \textbf{5-nitro-bufotenin} --- an autocatalytic
DMT analog found by promoting Phi from 𐑮 (complex-plane critical) to ⊙
(self-modeling critical):

\begin{verbatim}
DMT:      ⟨𐑨𐑸𐑩𐑗𐑞𐑘𐑔𐑵𐑮𐑫𐑕𐑭⟩   Phi=𐑮 (complex critical)
DMT-⊙:    ⟨𐑨𐑸𐑩𐑿𐑐𐑘𐑔𐑵⊙𐑫𐑕𐑭⟩   Phi=⊙ (self-modeling / autocatalytic)
\end{verbatim}

The molecule --- 5-nitro-bufotenin --- is a tryptamine with phenol
(EVALT), dimethylamine (EVALF), and nitro (ENGAGR), giving it all three
Dialetheia tokens simultaneously for autocatalytic template-directed
synthesis.

SMILES: \texttt{CN(C)CCC1=CNC2=CC=C(O)C({[}N+{]}(=O){[}O-{]})=C12}

\hypertarget{finder-tool}{%
\subsubsection{⊙-Finder Tool}\label{finder-tool}}

The standalone \textbf{⊙-Finder}
(\texttt{ig-docs/dmt-odot-discovery/odot\_finder.py}, 8.6 KB)
systematically finds ⊙-critical variants of any molecule:

\begin{verbatim}
# Find ⊙-critical variants of DMT
python3 /home/mrnob0dy666/imsgct/ig-docs/dmt-odot-discovery/odot_finder.py \
  --smiles "CN(C)CCC1=CNC2=CC=CC=C12" --name "DMT"

# Run demo mode (all test cases)
python3 /home/mrnob0dy666/imsgct/ig-docs/dmt-odot-discovery/odot_finder.py
\end{verbatim}

Results: \textbar{} Compound \textbar{} ⊙-Critical Candidates \textbar{}
\textbar----------\textbar----------------------\textbar{} \textbar{}
DMT \textbar{} 20 \textbar{} \textbar{} Serotonin \textbar{} 4
\textbar{} \textbar{} Bufotenin \textbar{} 4 \textbar{} \textbar{}
5-nitro-bufotenin \textbar{} Already ⊙-critical \textbar{}

\hypertarget{discovery-files}{%
\subsubsection{Discovery Files}\label{discovery-files}}

Located at \texttt{ig-docs/dmt-odot-discovery/}: - \texttt{README.md}
--- Full documentation (8.5 KB) - \texttt{5-nitro-bufotenin.cdxml} ---
ChemDraw drawing (4.9 KB) - \texttt{discovery.json} --- Structured
metadata (1.3 KB) - \texttt{dmt\_odot\_ob3ect.py} --- Self-verifying
ob3ect (1.9 KB) - \texttt{odot\_finder.py} --- ⊙-critical molecule
finder (8.6 KB) ---

\hypertarget{stress-testing}{%
\subsection{Stress Testing}\label{stress-testing}}

Two comprehensive stress test suites validate the entire platform across
edge cases, large inputs, boundary conditions, and cross-pipeline
integration. Both suites run to completion with 100\% pass rate.

\hypertarget{protein-genetics-34-tests-11-groups}{%
\subsubsection{Protein \& Genetics (34 tests, 11
groups)}\label{protein-genetics-34-tests-11-groups}}

\begin{verbatim}
python3 rebis.py run stress_test_proteins
# or directly:
python3 stress_test_proteins.py
\end{verbatim}

\begin{longtable}[]{@{}llll@{}}
\toprule
\begin{minipage}[b]{0.09\columnwidth}\raggedright
\#\strut
\end{minipage} & \begin{minipage}[b]{0.31\columnwidth}\raggedright
Test Group\strut
\end{minipage} & \begin{minipage}[b]{0.20\columnwidth}\raggedright
Tests\strut
\end{minipage} & \begin{minipage}[b]{0.29\columnwidth}\raggedright
Coverage\strut
\end{minipage}\tabularnewline
\midrule
\endhead
\begin{minipage}[t]{0.09\columnwidth}\raggedright
1\strut
\end{minipage} & \begin{minipage}[t]{0.31\columnwidth}\raggedright
Genetic Code\strut
\end{minipage} & \begin{minipage}[t]{0.20\columnwidth}\raggedright
4\strut
\end{minipage} & \begin{minipage}[t]{0.29\columnwidth}\raggedright
64-codon Frobenius-verified table, B4 invariant, lattice
operations\strut
\end{minipage}\tabularnewline
\begin{minipage}[t]{0.09\columnwidth}\raggedright
2\strut
\end{minipage} & \begin{minipage}[t]{0.31\columnwidth}\raggedright
Genetics B4\strut
\end{minipage} & \begin{minipage}[t]{0.20\columnwidth}\raggedright
3\strut
\end{minipage} & \begin{minipage}[t]{0.29\columnwidth}\raggedright
Nucleotide↔Belnap (T/B/F/N), structural distance, wobble pair
detection\strut
\end{minipage}\tabularnewline
\begin{minipage}[t]{0.09\columnwidth}\raggedright
3\strut
\end{minipage} & \begin{minipage}[t]{0.31\columnwidth}\raggedright
Genetic Tuples\strut
\end{minipage} & \begin{minipage}[t]{0.20\columnwidth}\raggedright
3\strut
\end{minipage} & \begin{minipage}[t]{0.29\columnwidth}\raggedright
7-stage tuple generation, structural distance DNA→quaternary, tier
consistency\strut
\end{minipage}\tabularnewline
\begin{minipage}[t]{0.09\columnwidth}\raggedright
4\strut
\end{minipage} & \begin{minipage}[t]{0.31\columnwidth}\raggedright
Genetic ASM\strut
\end{minipage} & \begin{minipage}[t]{0.20\columnwidth}\raggedright
1\strut
\end{minipage} & \begin{minipage}[t]{0.29\columnwidth}\raggedright
Paraconsistent abstract state machine program execution\strut
\end{minipage}\tabularnewline
\begin{minipage}[t]{0.09\columnwidth}\raggedright
5\strut
\end{minipage} & \begin{minipage}[t]{0.31\columnwidth}\raggedright
Gene→Protein Pipeline\strut
\end{minipage} & \begin{minipage}[t]{0.20\columnwidth}\raggedright
8\strut
\end{minipage} & \begin{minipage}[t]{0.29\columnwidth}\raggedright
Standard 7-stage, empty sequence, short (≤1 nt), start→stop, 100-AA,
multi-stop, no-stop, primitive activations\strut
\end{minipage}\tabularnewline
\begin{minipage}[t]{0.09\columnwidth}\raggedright
6\strut
\end{minipage} & \begin{minipage}[t]{0.31\columnwidth}\raggedright
SerpentRod\strut
\end{minipage} & \begin{minipage}[t]{0.20\columnwidth}\raggedright
5\strut
\end{minipage} & \begin{minipage}[t]{0.29\columnwidth}\raggedright
12-AA design, single-codon, empty, 200-AA, poly-Met homopolymer\strut
\end{minipage}\tabularnewline
\begin{minipage}[t]{0.09\columnwidth}\raggedright
7\strut
\end{minipage} & \begin{minipage}[t]{0.31\columnwidth}\raggedright
Ch3mpiler-SerpentRod\strut
\end{minipage} & \begin{minipage}[t]{0.20\columnwidth}\raggedright
3\strut
\end{minipage} & \begin{minipage}[t]{0.29\columnwidth}\raggedright
Caffeine (58-08-2), aspirin (50-78-2), cross-molecule
differentiation\strut
\end{minipage}\tabularnewline
\begin{minipage}[t]{0.09\columnwidth}\raggedright
8\strut
\end{minipage} & \begin{minipage}[t]{0.31\columnwidth}\raggedright
Cross-Integration\strut
\end{minipage} & \begin{minipage}[t]{0.20\columnwidth}\raggedright
2\strut
\end{minipage} & \begin{minipage}[t]{0.29\columnwidth}\raggedright
Gene→protein→fold, ch3mpiler→material bridge\strut
\end{minipage}\tabularnewline
\begin{minipage}[t]{0.09\columnwidth}\raggedright
9\strut
\end{minipage} & \begin{minipage}[t]{0.31\columnwidth}\raggedright
Antibody Designer\strut
\end{minipage} & \begin{minipage}[t]{0.20\columnwidth}\raggedright
1\strut
\end{minipage} & \begin{minipage}[t]{0.29\columnwidth}\raggedright
Epitope analysis → CDR → full VH antibody (53 AA, Frobenius ✓)\strut
\end{minipage}\tabularnewline
\begin{minipage}[t]{0.09\columnwidth}\raggedright
10\strut
\end{minipage} & \begin{minipage}[t]{0.31\columnwidth}\raggedright
PDB Validator\strut
\end{minipage} & \begin{minipage}[t]{0.20\columnwidth}\raggedright
1\strut
\end{minipage} & \begin{minipage}[t]{0.29\columnwidth}\raggedright
Structure validation, precision/recall metrics\strut
\end{minipage}\tabularnewline
\begin{minipage}[t]{0.09\columnwidth}\raggedright
11\strut
\end{minipage} & \begin{minipage}[t]{0.31\columnwidth}\raggedright
Edge Cases\strut
\end{minipage} & \begin{minipage}[t]{0.20\columnwidth}\raggedright
3\strut
\end{minipage} & \begin{minipage}[t]{0.29\columnwidth}\raggedright
All 3 stop codons (UAA/UAG/UGA), AUG start, non-standard nucleotides
(X/N/R/Y/?)\strut
\end{minipage}\tabularnewline
\bottomrule
\end{longtable}

\hypertarget{materials-26-tests-12-groups}{%
\subsubsection{Materials (26 tests, 12
groups)}\label{materials-26-tests-12-groups}}

\begin{verbatim}
python3 materials/stress_test_materials.py
\end{verbatim}

\begin{longtable}[]{@{}llll@{}}
\toprule
\begin{minipage}[b]{0.09\columnwidth}\raggedright
\#\strut
\end{minipage} & \begin{minipage}[b]{0.31\columnwidth}\raggedright
Test Group\strut
\end{minipage} & \begin{minipage}[b]{0.20\columnwidth}\raggedright
Tests\strut
\end{minipage} & \begin{minipage}[b]{0.29\columnwidth}\raggedright
Coverage\strut
\end{minipage}\tabularnewline
\midrule
\endhead
\begin{minipage}[t]{0.09\columnwidth}\raggedright
1\strut
\end{minipage} & \begin{minipage}[t]{0.31\columnwidth}\raggedright
IG Material Forge\strut
\end{minipage} & \begin{minipage}[t]{0.20\columnwidth}\raggedright
4\strut
\end{minipage} & \begin{minipage}[t]{0.29\columnwidth}\raggedright
All 8 predefined materials, empty name, report, list\strut
\end{minipage}\tabularnewline
\begin{minipage}[t]{0.09\columnwidth}\raggedright
2\strut
\end{minipage} & \begin{minipage}[t]{0.31\columnwidth}\raggedright
Sophick Forge\strut
\end{minipage} & \begin{minipage}[t]{0.20\columnwidth}\raggedright
4\strut
\end{minipage} & \begin{minipage}[t]{0.29\columnwidth}\raggedright
7-eagle cycle, progressive refinement, report, etch\strut
\end{minipage}\tabularnewline
\begin{minipage}[t]{0.09\columnwidth}\raggedright
3\strut
\end{minipage} & \begin{minipage}[t]{0.31\columnwidth}\raggedright
Frobenius Metamaterial\strut
\end{minipage} & \begin{minipage}[t]{0.20\columnwidth}\raggedright
3\strut
\end{minipage} & \begin{minipage}[t]{0.29\columnwidth}\raggedright
10×10 grid 20-cycle healing, load-heal cycle, export\strut
\end{minipage}\tabularnewline
\begin{minipage}[t]{0.09\columnwidth}\raggedright
4\strut
\end{minipage} & \begin{minipage}[t]{0.31\columnwidth}\raggedright
Thermal Rectifier\strut
\end{minipage} & \begin{minipage}[t]{0.20\columnwidth}\raggedright
2\strut
\end{minipage} & \begin{minipage}[t]{0.29\columnwidth}\raggedright
Diode asymmetry ratio, full simulation run\strut
\end{minipage}\tabularnewline
\begin{minipage}[t]{0.09\columnwidth}\raggedright
5\strut
\end{minipage} & \begin{minipage}[t]{0.31\columnwidth}\raggedright
Non-Qubit QC\strut
\end{minipage} & \begin{minipage}[t]{0.20\columnwidth}\raggedright
2\strut
\end{minipage} & \begin{minipage}[t]{0.29\columnwidth}\raggedright
all\_deltas computation, paradigm summary table\strut
\end{minipage}\tabularnewline
\begin{minipage}[t]{0.09\columnwidth}\raggedright
6\strut
\end{minipage} & \begin{minipage}[t]{0.31\columnwidth}\raggedright
Ouroboric Alloy\strut
\end{minipage} & \begin{minipage}[t]{0.20\columnwidth}\raggedright
3\strut
\end{minipage} & \begin{minipage}[t]{0.29\columnwidth}\raggedright
Init, mechanical test, heal cycle\strut
\end{minipage}\tabularnewline
\begin{minipage}[t]{0.09\columnwidth}\raggedright
7\strut
\end{minipage} & \begin{minipage}[t]{0.31\columnwidth}\raggedright
Critical Metamaterial\strut
\end{minipage} & \begin{minipage}[t]{0.20\columnwidth}\raggedright
2\strut
\end{minipage} & \begin{minipage}[t]{0.29\columnwidth}\raggedright
Susceptibility divergence, critical run\strut
\end{minipage}\tabularnewline
\begin{minipage}[t]{0.09\columnwidth}\raggedright
8\strut
\end{minipage} & \begin{minipage}[t]{0.31\columnwidth}\raggedright
Gap Closure\strut
\end{minipage} & \begin{minipage}[t]{0.20\columnwidth}\raggedright
1\strut
\end{minipage} & \begin{minipage}[t]{0.29\columnwidth}\raggedright
Enum completeness check\strut
\end{minipage}\tabularnewline
\begin{minipage}[t]{0.09\columnwidth}\raggedright
9\strut
\end{minipage} & \begin{minipage}[t]{0.31\columnwidth}\raggedright
Frobenius Exactor\strut
\end{minipage} & \begin{minipage}[t]{0.20\columnwidth}\raggedright
1\strut
\end{minipage} & \begin{minipage}[t]{0.29\columnwidth}\raggedright
All 4 designs (ω/τ/σ/ε)\strut
\end{minipage}\tabularnewline
\begin{minipage}[t]{0.09\columnwidth}\raggedright
10\strut
\end{minipage} & \begin{minipage}[t]{0.31\columnwidth}\raggedright
Materials Simulation\strut
\end{minipage} & \begin{minipage}[t]{0.20\columnwidth}\raggedright
1\strut
\end{minipage} & \begin{minipage}[t]{0.29\columnwidth}\raggedright
SelfHealingComposite + EternalMemorySim\strut
\end{minipage}\tabularnewline
\begin{minipage}[t]{0.09\columnwidth}\raggedright
11\strut
\end{minipage} & \begin{minipage}[t]{0.31\columnwidth}\raggedright
Frobenius Closure\strut
\end{minipage} & \begin{minipage}[t]{0.20\columnwidth}\raggedright
1\strut
\end{minipage} & \begin{minipage}[t]{0.29\columnwidth}\raggedright
ClosureDesign status\strut
\end{minipage}\tabularnewline
\begin{minipage}[t]{0.09\columnwidth}\raggedright
12\strut
\end{minipage} & \begin{minipage}[t]{0.31\columnwidth}\raggedright
Cross-Module\strut
\end{minipage} & \begin{minipage}[t]{0.20\columnwidth}\raggedright
2\strut
\end{minipage} & \begin{minipage}[t]{0.29\columnwidth}\raggedright
Forge→metamaterial, eagle→exactor\strut
\end{minipage}\tabularnewline
\bottomrule
\end{longtable}

\hypertarget{molecule-material-bridge}{%
\subsubsection{Molecule → Material
Bridge}\label{molecule-material-bridge}}

Derives material structural types directly from molecular input:

\begin{verbatim}
python3 materials/molecule_material_bridge.py --cas 58-08-2
python3 materials/molecule_material_bridge.py --smiles "CN1C=NC2=C1C(=O)N(C(=O)N2C)C"
\end{verbatim}

Produces: material tuple, tier assessment, Frobenius closure status,
catalog cross-reference.

\begin{center}\rule{0.5\linewidth}{0.5pt}\end{center}

\hypertarget{the-paraconsistent-kernel-rhr_p4rky}{%
\subsection{\texorpdfstring{The Paraconsistent Kernel
(\texttt{rhr\_p4rky/})}{The Paraconsistent Kernel (rhr\_p4rky/)}}\label{the-paraconsistent-kernel-rhr_p4rky}}

The paraconsistent kernel is a 32-module Python library providing Belnap
FOUR logic as universal substrate for computation that tolerates
contradiction. It was migrated from the standalone \texttt{p4rakernel/}
project and now lives as a subsystem within Red-Hot Rebis.

\hypertarget{architecture}{%
\subsubsection{Architecture}\label{architecture}}

\begin{longtable}[]{@{}lll@{}}
\toprule
\begin{minipage}[b]{0.23\columnwidth}\raggedright
Layer\strut
\end{minipage} & \begin{minipage}[b]{0.26\columnwidth}\raggedright
Module\strut
\end{minipage} & \begin{minipage}[b]{0.42\columnwidth}\raggedright
Description\strut
\end{minipage}\tabularnewline
\midrule
\endhead
\begin{minipage}[t]{0.23\columnwidth}\raggedright
\textbf{Logic}\strut
\end{minipage} & \begin{minipage}[t]{0.26\columnwidth}\raggedright
\texttt{belnap.py}\strut
\end{minipage} & \begin{minipage}[t]{0.42\columnwidth}\raggedright
4-valued Belnap logic (T/B/F/N) --- truth, both, false, neither\strut
\end{minipage}\tabularnewline
\begin{minipage}[t]{0.23\columnwidth}\raggedright
\textbf{Machine}\strut
\end{minipage} & \begin{minipage}[t]{0.26\columnwidth}\raggedright
\texttt{machine.py}\strut
\end{minipage} & \begin{minipage}[t]{0.42\columnwidth}\raggedright
Paraconsistent Abstract State Machine (ParaASM)\strut
\end{minipage}\tabularnewline
\begin{minipage}[t]{0.23\columnwidth}\raggedright
\textbf{Genetics}\strut
\end{minipage} & \begin{minipage}[t]{0.26\columnwidth}\raggedright
\texttt{genetic\_code.py}\strut
\end{minipage} & \begin{minipage}[t]{0.42\columnwidth}\raggedright
64-codon Frobenius-verified genetic code\strut
\end{minipage}\tabularnewline
\begin{minipage}[t]{0.23\columnwidth}\raggedright
\strut
\end{minipage} & \begin{minipage}[t]{0.26\columnwidth}\raggedright
\texttt{genetics\_b4.py}\strut
\end{minipage} & \begin{minipage}[t]{0.42\columnwidth}\raggedright
B4 lattice --- 64 codons, 7-stage tuple verification\strut
\end{minipage}\tabularnewline
\begin{minipage}[t]{0.23\columnwidth}\raggedright
\strut
\end{minipage} & \begin{minipage}[t]{0.26\columnwidth}\raggedright
\texttt{genetic\_tuples.py}\strut
\end{minipage} & \begin{minipage}[t]{0.42\columnwidth}\raggedright
Tuple encodings for genes/codons/proteins\strut
\end{minipage}\tabularnewline
\begin{minipage}[t]{0.23\columnwidth}\raggedright
\strut
\end{minipage} & \begin{minipage}[t]{0.26\columnwidth}\raggedright
\texttt{genetic\_asm.py}\strut
\end{minipage} & \begin{minipage}[t]{0.42\columnwidth}\raggedright
Genetic abstract state machine\strut
\end{minipage}\tabularnewline
\begin{minipage}[t]{0.23\columnwidth}\raggedright
\textbf{Pipeline}\strut
\end{minipage} & \begin{minipage}[t]{0.26\columnwidth}\raggedright
\texttt{gene\_to\_protein\_pipeline.py}\strut
\end{minipage} & \begin{minipage}[t]{0.42\columnwidth}\raggedright
Full gene→protein translation (7-stage B4)\strut
\end{minipage}\tabularnewline
\begin{minipage}[t]{0.23\columnwidth}\raggedright
\strut
\end{minipage} & \begin{minipage}[t]{0.26\columnwidth}\raggedright
\texttt{demo\_gene\_to\_protein.py}\strut
\end{minipage} & \begin{minipage}[t]{0.42\columnwidth}\raggedright
Pipeline demonstration\strut
\end{minipage}\tabularnewline
\begin{minipage}[t]{0.23\columnwidth}\raggedright
\strut
\end{minipage} & \begin{minipage}[t]{0.26\columnwidth}\raggedright
\texttt{run\_gene\_pipeline.py}\strut
\end{minipage} & \begin{minipage}[t]{0.42\columnwidth}\raggedright
Pipeline CLI runner\strut
\end{minipage}\tabularnewline
\begin{minipage}[t]{0.23\columnwidth}\raggedright
\textbf{Proteins}\strut
\end{minipage} & \begin{minipage}[t]{0.26\columnwidth}\raggedright
\texttt{serpent\_rod.py}\strut
\end{minipage} & \begin{minipage}[t]{0.42\columnwidth}\raggedright
Foundational Frobenius morphism RNA→\{sequence+fold\}; base
\texttt{SerpentRod} class\strut
\end{minipage}\tabularnewline
\begin{minipage}[t]{0.23\columnwidth}\raggedright
\strut
\end{minipage} & \begin{minipage}[t]{0.26\columnwidth}\raggedright
\texttt{serpent\_rod\_v2.py}\strut
\end{minipage} & \begin{minipage}[t]{0.42\columnwidth}\raggedright
\texttt{SerpentRodV2}: 3D backbone from φ/ψ angles, geometry-based
contacts, energy scoring\strut
\end{minipage}\tabularnewline
\begin{minipage}[t]{0.23\columnwidth}\raggedright
\textbf{Antibodies}\strut
\end{minipage} & \begin{minipage}[t]{0.26\columnwidth}\raggedright
\texttt{antibody\_designer.py}\strut
\end{minipage} & \begin{minipage}[t]{0.42\columnwidth}\raggedright
Computational antibody design from IG types\strut
\end{minipage}\tabularnewline
\begin{minipage}[t]{0.23\columnwidth}\raggedright
\textbf{Validation}\strut
\end{minipage} & \begin{minipage}[t]{0.26\columnwidth}\raggedright
\texttt{pdb\_validator.py}\strut
\end{minipage} & \begin{minipage}[t]{0.42\columnwidth}\raggedright
PDB structure validation against IG types; \texttt{extract\_sequence()}
round-trips SEQRES\strut
\end{minipage}\tabularnewline
\begin{minipage}[t]{0.23\columnwidth}\raggedright
\textbf{Bridges}\strut
\end{minipage} & \begin{minipage}[t]{0.26\columnwidth}\raggedright
\texttt{ch3mpiler\_bridge.py}\strut
\end{minipage} & \begin{minipage}[t]{0.42\columnwidth}\raggedright
Ch3mpiler ↔ kernel bridge\strut
\end{minipage}\tabularnewline
\begin{minipage}[t]{0.23\columnwidth}\raggedright
\strut
\end{minipage} & \begin{minipage}[t]{0.26\columnwidth}\raggedright
\texttt{ch3mpiler\_ob3ect\_bridge.py}\strut
\end{minipage} & \begin{minipage}[t]{0.42\columnwidth}\raggedright
Ch3mpiler ↔ ob3ect bridge\strut
\end{minipage}\tabularnewline
\begin{minipage}[t]{0.23\columnwidth}\raggedright
\strut
\end{minipage} & \begin{minipage}[t]{0.26\columnwidth}\raggedright
\texttt{ch3mpiler\_serpentrod\_pipeline.py}\strut
\end{minipage} & \begin{minipage}[t]{0.42\columnwidth}\raggedright
\textbf{v4} --- Weighted bond-preserving fusion, word-boundary FG
matching, specificity-ranked disconnections, Frobenius-exact complement
mapping, CAS lookup, molecule-specific catalytic sites\strut
\end{minipage}\tabularnewline
\begin{minipage}[t]{0.23\columnwidth}\raggedright
\textbf{Physics}\strut
\end{minipage} & \begin{minipage}[t]{0.26\columnwidth}\raggedright
\texttt{hadron\_belnap.py}\strut
\end{minipage} & \begin{minipage}[t]{0.42\columnwidth}\raggedright
Hadronic Belnap-state classification\strut
\end{minipage}\tabularnewline
\begin{minipage}[t]{0.23\columnwidth}\raggedright
\strut
\end{minipage} & \begin{minipage}[t]{0.26\columnwidth}\raggedright
\texttt{exotic\_hadron\_belnap.py}\strut
\end{minipage} & \begin{minipage}[t]{0.42\columnwidth}\raggedright
Exotic hadron Belnap classification\strut
\end{minipage}\tabularnewline
\begin{minipage}[t]{0.23\columnwidth}\raggedright
\strut
\end{minipage} & \begin{minipage}[t]{0.26\columnwidth}\raggedright
\texttt{quark\_belnap.py}\strut
\end{minipage} & \begin{minipage}[t]{0.42\columnwidth}\raggedright
Quark Belnap-state analysis\strut
\end{minipage}\tabularnewline
\begin{minipage}[t]{0.23\columnwidth}\raggedright
\strut
\end{minipage} & \begin{minipage}[t]{0.26\columnwidth}\raggedright
\texttt{orbital\_belnap.py}\strut
\end{minipage} & \begin{minipage}[t]{0.42\columnwidth}\raggedright
Orbital Belnap-state analysis\strut
\end{minipage}\tabularnewline
\begin{minipage}[t]{0.23\columnwidth}\raggedright
\textbf{Analysis}\strut
\end{minipage} & \begin{minipage}[t]{0.26\columnwidth}\raggedright
\texttt{clu\_power\_law.py}\strut
\end{minipage} & \begin{minipage}[t]{0.42\columnwidth}\raggedright
Clustering power-law analysis\strut
\end{minipage}\tabularnewline
\begin{minipage}[t]{0.23\columnwidth}\raggedright
\strut
\end{minipage} & \begin{minipage}[t]{0.26\columnwidth}\raggedright
\texttt{frobenius\_filtration.py}\strut
\end{minipage} & \begin{minipage}[t]{0.42\columnwidth}\raggedright
Frobenius-verified filtration\strut
\end{minipage}\tabularnewline
\begin{minipage}[t]{0.23\columnwidth}\raggedright
\textbf{New modules}\strut
\end{minipage} & \begin{minipage}[t]{0.26\columnwidth}\raggedright
\strut
\end{minipage} & \begin{minipage}[t]{0.42\columnwidth}\raggedright
\strut
\end{minipage}\tabularnewline
\begin{minipage}[t]{0.23\columnwidth}\raggedright
\texttt{decay\_chain.py}\strut
\end{minipage} & \begin{minipage}[t]{0.26\columnwidth}\raggedright
Nuclear decay as IMASM winding --- 5 natural series to Frobenius fixed
point\strut
\end{minipage} & \begin{minipage}[t]{0.42\columnwidth}\raggedright
\strut
\end{minipage}\tabularnewline
\begin{minipage}[t]{0.23\columnwidth}\raggedright
\texttt{belnap\_c4.py}\strut
\end{minipage} & \begin{minipage}[t]{0.26\columnwidth}\raggedright
Belnap C4 logic variant --- contradiction-majority lattice\strut
\end{minipage} & \begin{minipage}[t]{0.42\columnwidth}\raggedright
\strut
\end{minipage}\tabularnewline
\begin{minipage}[t]{0.23\columnwidth}\raggedright
\textbf{papers/}\strut
\end{minipage} & \begin{minipage}[t]{0.26\columnwidth}\raggedright
3 millennium documents\strut
\end{minipage} & \begin{minipage}[t]{0.42\columnwidth}\raggedright
\strut
\end{minipage}\tabularnewline
\begin{minipage}[t]{0.23\columnwidth}\raggedright
\texttt{all\_millennium\_solved.md}\strut
\end{minipage} & \begin{minipage}[t]{0.26\columnwidth}\raggedright
Unified Millennium solution framework\strut
\end{minipage} & \begin{minipage}[t]{0.42\columnwidth}\raggedright
\strut
\end{minipage}\tabularnewline
\begin{minipage}[t]{0.23\columnwidth}\raggedright
\texttt{belnap\_qm.md}\strut
\end{minipage} & \begin{minipage}[t]{0.26\columnwidth}\raggedright
Belnap quantum mechanics formalization\strut
\end{minipage} & \begin{minipage}[t]{0.42\columnwidth}\raggedright
\strut
\end{minipage}\tabularnewline
\begin{minipage}[t]{0.23\columnwidth}\raggedright
\texttt{millennium\_barriers.md}\strut
\end{minipage} & \begin{minipage}[t]{0.26\columnwidth}\raggedright
Barrier taxonomy across all 7 Millennium Problems\strut
\end{minipage} & \begin{minipage}[t]{0.42\columnwidth}\raggedright
\strut
\end{minipage}\tabularnewline
\begin{minipage}[t]{0.23\columnwidth}\raggedright
\textbf{Utilities}\strut
\end{minipage} & \begin{minipage}[t]{0.26\columnwidth}\raggedright
\texttt{kernel.py}\strut
\end{minipage} & \begin{minipage}[t]{0.42\columnwidth}\raggedright
Kernel core initialization\strut
\end{minipage}\tabularnewline
\begin{minipage}[t]{0.23\columnwidth}\raggedright
\strut
\end{minipage} & \begin{minipage}[t]{0.26\columnwidth}\raggedright
\texttt{pipeline\_fix.py}\strut
\end{minipage} & \begin{minipage}[t]{0.42\columnwidth}\raggedright
Pipeline repair utilities\strut
\end{minipage}\tabularnewline
\bottomrule
\end{longtable}

\hypertarget{key-commands}{%
\subsubsection{Key commands}\label{key-commands}}

\begin{verbatim}
# Full genetics test suite (B4 lattice, 64-codon, gene→protein, Phi gate, ParaASM)
python3 test_genetics.py

# Comprehensive protein & genetics stress test (34 tests, 11 groups)
python3 rebis.py run stress_test_proteins

# Single gene pipeline
python3 rebis.py run run_gene_pipeline --gene INS

# Genetic code analysis
python3 rebis.py run genetic_code

# Antibody design from viral epitopes
python3 rebis.py run antibody_designer

# PDB validation
python3 rebis.py run pdb_validator

# Hadron Belnap classification
python3 rebis.py run hadron_belnap
python3 rebis.py run exotic_hadron_belnap
python3 rebis.py run quark_belnap
\end{verbatim}

\begin{center}\rule{0.5\linewidth}{0.5pt}\end{center}

\hypertarget{gene-imscriber-gene_imscriber}{%
\subsection{\texorpdfstring{Gene Imscriber
(\texttt{gene\_imscriber/})}{Gene Imscriber (gene\_imscriber/)}}\label{gene-imscriber-gene_imscriber}}

The Gene Imscriber provides IG-native genetic compilation --- structural
types directly to codon optimization, CRISPR guide design, chimera
construction, and Frobenius-verified base/prime editing.

\begin{longtable}[]{@{}ll@{}}
\toprule
Module & Description\tabularnewline
\midrule
\endhead
\texttt{engine.py} & Core engine --- structural type → codon
optimization\tabularnewline
\texttt{genetics\_ig\_prelim.py} & Preliminary IG-to-genetics
mapping\tabularnewline
\texttt{genetics\_ig\_promotions.py} & IG promotion paths for
genetics\tabularnewline
\texttt{genetics\_qs.py} & Quantum superposition in genetic
space\tabularnewline
\texttt{ig\_genetics\_answer.py} & IG-native genetics
answers\tabularnewline
\texttt{tuples.py} & Genetic tuple definitions\tabularnewline
\bottomrule
\end{longtable}

The Gene Imscriber also contains \texttt{scripts/} with GUIDE-seq
analysis, base editor stratum analysis, clinical safety pipelines, and
SRA GUIDE-seq processing.

\begin{verbatim}
python3 rebis.py run gene --help
\end{verbatim}

\begin{center}\rule{0.5\linewidth}{0.5pt}\end{center}

\hypertarget{popular-protein-analysis-popular_protein}{%
\subsection{\texorpdfstring{Popular Protein Analysis
(\texttt{popular\_protein/})}{Popular Protein Analysis (popular\_protein/)}}\label{popular-protein-analysis-popular_protein}}

A comparison toolkit for validating SerpentRod predictions against
crystal structures:

\begin{verbatim}
python3 rebis.py run compare_exact          # Exact φ/ψ angle comparison
python3 rebis.py run compare_structures     # Structure-level comparison
python3 rebis.py run comprehensive_comparison  # Full multi-metric report
\end{verbatim}

Reference PDB files are in \texttt{pdb/}: 1L2Y, 1UBQ, 1VII, 1ZDD.

\begin{center}\rule{0.5\linewidth}{0.5pt}\end{center}

\hypertarget{key-concepts}{%
\subsection{Key Concepts}\label{key-concepts}}

\hypertarget{ig-tuples}{%
\subsubsection{IG Tuples}\label{ig-tuples}}

Every structural entity carries a 12-primitive tuple. Each position is a
primitive with a value from its ordinal set. Tuples are the sole
carriers of structural information --- no external parameters, no
assumed constants. The catalog (\texttt{shared/IG\_catalog.json})
contains \textbf{4,027+} verified entries (up from 3,297).

\hypertarget{the-frobenius-condition}{%
\subsubsection{The Frobenius Condition}\label{the-frobenius-condition}}

\(\mu\circ\delta=\text{id}\). When this holds, the structural type is
self-consistent --- comultiplication followed by multiplication returns
the identity. A Frobenius ✅ means the layer's tuple is internally
coherent. A ❌ means there is a primitive conflict to resolve.

\hypertarget{tiers-ouroboricity}{%
\subsubsection{Tiers (Ouroboricity)}\label{tiers-ouroboricity}}

Structural types are classified into ouroboricity tiers based on their
\(\Phi\) (parity/symmetry) and \(\Omega\) (winding) primitives:

\begin{longtable}[]{@{}lll@{}}
\toprule
\begin{minipage}[b]{0.21\columnwidth}\raggedright
Tier\strut
\end{minipage} & \begin{minipage}[b]{0.39\columnwidth}\raggedright
Condition\strut
\end{minipage} & \begin{minipage}[b]{0.32\columnwidth}\raggedright
Meaning\strut
\end{minipage}\tabularnewline
\midrule
\endhead
\begin{minipage}[t]{0.21\columnwidth}\raggedright
\(O₀\)\strut
\end{minipage} & \begin{minipage}[t]{0.39\columnwidth}\raggedright
\(\Phi \leq 𐑬\), \(\Omega = 𐑷\)\strut
\end{minipage} & \begin{minipage}[t]{0.32\columnwidth}\raggedright
No winding, no self-modeling --- trivial\strut
\end{minipage}\tabularnewline
\begin{minipage}[t]{0.21\columnwidth}\raggedright
\(O₁\)\strut
\end{minipage} & \begin{minipage}[t]{0.39\columnwidth}\raggedright
\(\Phi \leq 𐑿\), \(\Omega \geq 𐑴\)\strut
\end{minipage} & \begin{minipage}[t]{0.32\columnwidth}\raggedright
Quantum parity, \(\mathbb{Z}_2\) winding --- simple loop\strut
\end{minipage}\tabularnewline
\begin{minipage}[t]{0.21\columnwidth}\raggedright
\(O₂\)\strut
\end{minipage} & \begin{minipage}[t]{0.39\columnwidth}\raggedright
\(\Phi \geq 𐑬\), \(\Omega \geq 𐑭\)\strut
\end{minipage} & \begin{minipage}[t]{0.32\columnwidth}\raggedright
Partial symmetry, \(\mathbb{Z}\) winding --- structured loop\strut
\end{minipage}\tabularnewline
\begin{minipage}[t]{0.21\columnwidth}\raggedright
\(O_∞\)\strut
\end{minipage} & \begin{minipage}[t]{0.39\columnwidth}\raggedright
\(\Phi = 𐑹\), \(\Omega \geq 𐑟\), \(\odot\)=⊙\strut
\end{minipage} & \begin{minipage}[t]{0.32\columnwidth}\raggedright
Frobenius-special, non-Abelian winding, self-modeling gate open\strut
\end{minipage}\tabularnewline
\bottomrule
\end{longtable}

\hypertarget{imasm-token-system}{%
\subsubsection{IMASM Token System}\label{imasm-token-system}}

The IMASM token algebra has 12 tokens organized in semantic families:

\begin{longtable}[]{@{}lll@{}}
\toprule
Token & Category & Chemical Meaning\tabularnewline
\midrule
\endhead
VINIT & Scaffold & Carbon backbone, ring system\tabularnewline
AFWD & Forward reactivity & Electrophilic site, carbonyl\tabularnewline
AREV & Reverse reactivity & Nucleophilic site,
alcohol/amine\tabularnewline
FFUSE & Fusion site & Bond formation site\tabularnewline
FSPLIT & Disconnection point & Retrosynthetic cut site\tabularnewline
EVALT & Acidic & Carboxylic acid, phenol\tabularnewline
EVALF & Basic & Amine, amidine\tabularnewline
ENGAGR & Ambident & Nitro, nitrile --- multi-site
reactivity\tabularnewline
CLINK & Linker & Conjugated system, spacer\tabularnewline
TANCH & Terminal anchor & Terminal functional group\tabularnewline
IFIX & Fixation & Cyclization, ring closure\tabularnewline
IMSCRIB & Self-imscribing & Aromatic ring, self-consistent
system\tabularnewline
\bottomrule
\end{longtable}

\hypertarget{dialetheia-triad}{%
\subsubsection{Dialetheia Triad}\label{dialetheia-triad}}

A compound is \textbf{Dialetheia-complete} when it simultaneously
contains three IMASM tokens: - \textbf{EVALT} (acidic functional group
--- carboxylic acid, phenol) - \textbf{EVALF} (basic functional group
--- amine, amidine) - \textbf{ENGAGR} (ambident functional group ---
nitro, nitrile)

This triad enables autocatalytic template-directed synthesis --- a
necessary condition for ⊙ criticality in molecules.

\hypertarget{cross-domain-analogies}{%
\subsubsection{Cross-Domain Analogies}\label{cross-domain-analogies}}

\hypertarget{because-all-4027-catalog-entries-share-the-same-12-primitive-typing-a-compounds-ig-type-can-be-compared-to-any-other-system-materials-languages-consciousness-states-mathematical-theorems.-this-is-what-makes-the-imasm-compound-pipeline-powerful-a-molecules-structural-fingerprint-reveals-kinship-with-systems-that-no-conventional-chemical-descriptor-could-reach.}{%
\subsection{Because all 4,027+ catalog entries share the same
12-primitive typing, a compound's IG type can be compared to any other
system --- materials, languages, consciousness states, mathematical
theorems. This is what makes the IMASM compound pipeline powerful: a
molecule's structural fingerprint reveals kinship with systems that no
conventional chemical descriptor could
reach.}\label{because-all-4027-catalog-entries-share-the-same-12-primitive-typing-a-compounds-ig-type-can-be-compared-to-any-other-system-materials-languages-consciousness-states-mathematical-theorems.-this-is-what-makes-the-imasm-compound-pipeline-powerful-a-molecules-structural-fingerprint-reveals-kinship-with-systems-that-no-conventional-chemical-descriptor-could-reach.}}

\hypertarget{common-workflows}{%
\subsection{Common Workflows}\label{common-workflows}}

\hypertarget{design-a-protein-from-scratch}{%
\subsubsection{Design a protein from
scratch}\label{design-a-protein-from-scratch}}

\begin{verbatim}
# 1. Check where proteins live in the CLINK chain
python3 rebis.py clink layer cell

# 2. Run SerpentRod on built-in test proteins
python3 rebis.py run serpentrod

# 3. Generate actionable organism package (includes protein.fasta + protein_coords.pdb)
python3 rebis.py pipeline actionable
\end{verbatim}

\hypertarget{design-a-catalytic-site-for-a-molecule}{%
\subsubsection{Design a catalytic site for a
molecule}\label{design-a-catalytic-site-for-a-molecule}}

\begin{verbatim}
# 1. CAS lookup (recommended — resolves molecule identity)
# 1. Direct retrosynthesis with SMILES input (recommended — exact)
python3 rebis.py run ch3mpiler --smiles "CN1C=NC2=C1C(=O)N(C(=O)N2C)C" --retrosynthesis

# 2. CAS lookup (also resolves SMILES)
python3 rebis.py run ch3mpiler_serpentrod_pipeline --cas 58-08-2

# 3. Direct molecule name (may not resolve if name is synthetic)
python3 rebis.py run ch3mpiler_serpentrod_pipeline --target "ibuprofen"

# 3. Retrosynthetic pair
python3 rebis.py run ch3mpiler_serpentrod_pipeline --start "salicylic_acid" --target "aspirin"

# 4. JSON output for programmatic use
python3 rebis.py run ch3mpiler_serpentrod_pipeline --cas 58-08-2 --json
\end{verbatim}

\hypertarget{encode-a-molecule-as-an-imasm-structural-type}{%
\subsubsection{Encode a molecule as an IMASM structural
type}\label{encode-a-molecule-as-an-imasm-structural-type}}

\begin{verbatim}
# 1. Encode a single compound
python3 rebis.py imas compound --smiles "CC(=O)OC1=CC=CC=C1C(=O)O"

# 2. Find cross-domain analogies
python3 rebis.py imas analogies --smiles "CN(C)CCC1=CNC2=CC=CC=C12" --limit 15

# 3. Analyze a reaction
python3 rebis.py imas reaction --smiles "CC(=O)C1=CC=CC=C1" --product "CC(C1=CC=CC=C1)(C)O"

# 4. Register a new compound for future lookups
python3 rebis.py imas register --smiles "CC(C)C1=CC=C(C=C1)C(C)C(=O)O" --name "ibuprofen"

# 5. Browse the compound catalog (static)
less INDEX.md    # §4 — Compound IMASM Catalog
\end{verbatim}

\hypertarget{find--critical-autocatalytic-molecules}{%
\subsubsection{Find ⊙-critical autocatalytic
molecules}\label{find--critical-autocatalytic-molecules}}

\begin{verbatim}
# 1. Run the ⊙-Finder on DMT
python3 /home/mrnob0dy666/imsgct/ig-docs/dmt-odot-discovery/odot_finder.py \
  --smiles "CN(C)CCC1=CNC2=CC=CC=C12" --name "DMT"

# 2. Run in demo mode (all test cases)
python3 /home/mrnob0dy666/imsgct/ig-docs/dmt-odot-discovery/odot_finder.py

# 3. Explore the discovery documentation
less /home/mrnob0dy666/imsgct/ig-docs/dmt-odot-discovery/README.md
\end{verbatim}

\hypertarget{work-with-the-paraconsistent-kernel}{%
\subsubsection{Work with the paraconsistent
kernel}\label{work-with-the-paraconsistent-kernel}}

\begin{verbatim}
# Run the full genetics test suite
python3 test_genetics.py

# Comprehensive stress test (34 tests across 11 groups)
python3 rebis.py run stress_test_proteins

# Translate a gene through the 7-stage B4 pipeline
python3 rebis.py run run_gene_pipeline --gene INS

# Analyze the 64-codon Frobenius-verified genetic code
python3 rebis.py run genetic_code

# Design a serpent rod protein
python3 rebis.py run serpent_rod_v2
\end{verbatim}

\hypertarget{investigate-a-material-type}{%
\subsubsection{Investigate a material
type}\label{investigate-a-material-type}}

\begin{verbatim}
# 1. Browse the materials catalog (static reference)
less INDEX.md                          # §3 — 8 predefined materials

# 2. Run a dynamic simulation
python3 rebis.py materials frobenius   # Frobenius closure verification
python3 rebis.py materials ouroboric   # Crack healing simulation

# 3. Eagle Cycle for O_∞ substrate preparation
python3 rebis.py materials sophick --name eagle_9_sophick

# 4. Forge the material design file
python3 rebis.py materials forge --name frobenius_composite

# 5. Run comprehensive materials stress test
python3 materials/stress_test_materials.py

# 6. Derive material type from molecule
python3 materials/molecule_material_bridge.py --cas 58-08-2
\end{verbatim}

\hypertarget{map-an-imasm-sequence-to-clink}{%
\subsubsection{Map an IMASM sequence to
CLINK}\label{map-an-imasm-sequence-to-clink}}

\begin{verbatim}
# 1. Browse the 12 IMASM canonicals (static reference)
less INDEX.md                          # §2 — 12 canonicals with tuples and tiers

# 2. Inspect a specific canonical's bridge to CLINK (dynamic)
python3 rebis.py imas bridge --canonical I_Dialetheic_Bootstrap

# 3. Compute activation energy to organism layer
python3 rebis.py imas energy --canonical I_Dialetheic_Bootstrap --layer L8_Organism

# 4. Monte Carlo Frobenius density estimation
python3 rebis.py imas hunt --samples 100000
\end{verbatim}

\hypertarget{analyze-viral-epitopes-for-antibody-design}{%
\subsubsection{Analyze viral epitopes for antibody
design}\label{analyze-viral-epitopes-for-antibody-design}}

\begin{verbatim}
python3 rebis.py run run_antibody
\end{verbatim}

Produces CDR sequences and structural complementarity scores for each
epitope target.

\hypertarget{mitochondrial-gene-analysis}{%
\subsubsection{Mitochondrial gene
analysis}\label{mitochondrial-gene-analysis}}

\begin{verbatim}
python3 rebis.py run mito_pipeline
\end{verbatim}

Processes all 13 MT genes, reports primitive activations (up to 9/12 per
gene), Frobenius status, and 7-stage B4 pipeline per gene (dna\_gene →
pre\_mrna → mature\_mrna → nascent\_polypeptide → secondary\_structure →
tertiary\_structure → quaternary\_structure).

\hypertarget{validate-serpentrod-predictions-against-crystal-structures}{%
\subsubsection{Validate SerpentRod predictions against crystal
structures}\label{validate-serpentrod-predictions-against-crystal-structures}}

\begin{verbatim}
python3 rebis.py run comprehensive_comparison
python3 rebis.py run compare_exact
\end{verbatim}

\hypertarget{hadron-particle-physics-classification}{%
\subsubsection{Hadron / particle physics
classification}\label{hadron-particle-physics-classification}}

\begin{verbatim}
python3 rebis.py run hadron_belnap
python3 rebis.py run exotic_hadron_belnap
python3 rebis.py run quark_belnap
\end{verbatim}

\begin{center}\rule{0.5\linewidth}{0.5pt}\end{center}

\hypertarget{troubleshooting}{%
\subsection{Troubleshooting}\label{troubleshooting}}

\textbf{\texttt{verify} shows ❌ for a module} Run
\texttt{python3\ -c\ "from\ \textless{}module\textgreater{}\ import\ *"}
directly to see the import error. Most common cause: a dependency
package not installed. Fix:

\begin{verbatim}
uv pip install -e /home/mrnob0dy666/imsgct/omonad_OS
uv pip install -e /home/mrnob0dy666/imsgct/imasmic_core
\end{verbatim}

\textbf{\texttt{omonad\_bridge} shows
\texttt{omonad\_available:\ false}}

\begin{verbatim}
uv pip install -e /home/mrnob0dy666/imsgct/omonad_OS
uv pip install -e /home/mrnob0dy666/imsgct/imasmic_core
\end{verbatim}

\textbf{\texttt{pipeline\ bridges} shows ❌ for \texttt{biology\_sim} or
\texttt{ouroboric\_telomere}} These two bridges are structurally
excluded --- modules exist on disk but are not wired to the bridge API.
Not a crash condition; the pipeline runs without them.

\textbf{\texttt{imas\ compound} fails with \texttt{ModuleNotFoundError}}
Ensure the package is installed:

\begin{verbatim}
uv pip install -e /home/mrnob0dy666/imsgct/red-hot_rebis
\end{verbatim}

The compound pipeline imports from \texttt{imas.compound\_imasm},
\texttt{imas.arranger}, and \texttt{imas.ig\_bridge} --- all part of the
Rebis package.

\textbf{\texttt{imas\ analogies} returns fewer results than
\texttt{-\/-limit}} The analogies search filters by structural distance.
If fewer neighbors exist within d ≤ 5, the result set will be smaller.
Try a higher limit. The compound must first be encodable (valid SMILES
with detectable functional groups).

\textbf{\texttt{imas\ reaction} misclassifies a reaction} The reaction
classifier uses token delta analysis (AFWD→AREV swap for reductions,
etc.). Unusual reaction mechanisms may fall outside the 7 built-in
classes. Check the raw token delta for structural insight.

\textbf{RDKit import errors (crystal designer, odot\_finder)} RDKit is
optional --- needed only for \texttt{molecular\_crystal\_designer.py}
and \texttt{odot\_finder.py}. Install with:

\begin{verbatim}
uv pip install rdkit-pypi
\end{verbatim}

\textbf{SerpentRod accuracy below 100\%} Fragment naming accuracy
(83--88\%) reflects genuine structural ambiguity in cleavage products.
SP accuracy should be 100\% on the standard test set; below that
indicates non-canonical signal peptide architecture in the test
sequence.

\textbf{Ch3mpiler retrosynthesis shows \texttt{C1=CC=C(C=C1)CC(C)C(=O)O}
for all precursors} This is the old placeholder SMILES --- the target
name couldn't be resolved. Use
\texttt{-\/-smiles\ "\textless{}exact\ SMILES\textgreater{}"} to provide
the molecular structure directly. The
\texttt{bond\_fragment\_integrator.py} will then cut real bonds and
display real fragment SMILES. If the target name does resolve
(e.g.~\texttt{-\/-target\ aspirin}), SMILES are auto-resolved via
OPSIN/PubChem. A bold red error message now appears when name resolution
fails.

\textbf{Ch3mpiler-SerpentRod produces identical sites for different
molecules} This was the pre-v4 MAX fusion bug. Update to the latest
pipeline:

\begin{verbatim}
cd /home/mrnob0dy666/imsgct/red-hot_rebis
git pull  # or ensure rhr_p4rky/ch3mpiler_serpentrod_pipeline.py has v4 weighted fusion
\end{verbatim}

Verify differentiation with:

\begin{verbatim}
python3 rebis.py run ch3mpiler_serpentrod_pipeline --cas 58-08-2  # caffeine
python3 rebis.py run ch3mpiler_serpentrod_pipeline --cas 17699-14-8  # tricyclic sesquiterpene
\end{verbatim}

Different AA sequences and catalytic triads confirm the fix.

\textbf{Genetics modules crash on unknown nucleotides} Fixed in v2.3.0
--- \texttt{genetics\_b4.py} \texttt{nucleotide\_to\_belnap} returns
\texttt{Belnap.B} for X/N/R/Y/?. \texttt{gene\_to\_protein\_pipeline.py}
handles empty/short sequences gracefully.

\textbf{\texttt{clink\ layer} gives ``No layer matching\ldots{}''} Use a
substring of the official layer name from INDEX.md §1 or
\texttt{python3\ rebis.py\ clink\ layer\ -\/-help}. Examples:
\texttt{quark} → \texttt{Frustrated\ Belnap5\ (Quarks)}, \texttt{tissue}
→ \texttt{Tissue/Organ}, \texttt{meiosis} → \texttt{Meiosis\ (Gametes)}.

\textbf{\texttt{materials\ forge} gives usage without generating} Always
specify \texttt{-\/-name\ \textless{}material\textgreater{}} or
\texttt{-\/-all}. Running without a name prints the usage message.

\textbf{\texttt{materials\ sophick} or \texttt{materials\ exactor} shows
usage prompt} These now require \texttt{-\/-name}. Run with
\texttt{-\/-help} to see valid name options, or browse INDEX.md §3.

\textbf{\texttt{run} target not found} Use
\texttt{python3\ rebis.py\ run\ list} to see all 39 discoverable
targets. Some targets require the package to be installed
(\texttt{uv\ pip\ install\ -e\ .}).

\textbf{\texttt{ModuleNotFoundError:\ No\ module\ named\ \textquotesingle{}rhr\_p4rky\textquotesingle{}}}
The rhr\_p4rky package must be installed:

\begin{verbatim}
uv pip install -e /home/mrnob0dy666/imsgct/red-hot_rebis
\end{verbatim}

\textbf{Paraconsistent kernel tests fail} Ensure the B4 lattice modules
are importable:

\begin{verbatim}
python3 -c "from rhr_p4rky.belnap import BelnapValue; print('OK')"
\end{verbatim}

\textbf{Stress tests fail} Run directly to see detailed tracebacks:

\begin{verbatim}
python3 stress_test_proteins.py          # Protein & genetics (34 tests)
python3 materials/stress_test_materials.py  # Materials (26 tests)
\end{verbatim}

Individual test failures report the exact module, function, and
exception.

\textbf{Where did \texttt{clink\ report}, \texttt{clink\ list},
\texttt{imas\ report}, \texttt{materials\ list},
\texttt{materials\ report} go?} These static-data commands were removed
from \texttt{rebis.py}. All their data now lives in \textbf{INDEX.md}
--- a plain-text browsable reference. Every removed command's
\texttt{-\/-help} points to the relevant INDEX.md section.

\begin{center}\rule{0.5\linewidth}{0.5pt}\end{center}

\begin{center}\rule{0.5\linewidth}{0.5pt}\end{center}

\hypertarget{at-ars-therapeutica}{%
\subsubsection{\texorpdfstring{\texttt{at} ---
Ars\,Therapeutica}{at --- Ars\,Therapeutica}}\label{at-ars-therapeutica}}

Grammar-derived optimal therapy design and structural diagnosis. Maps
diseases onto the 12-primitive crystal, computes structural distance to
healthy type, and generates the promotion path (therapy protocol) to
close the gap.

\begin{verbatim}
python3 rebis.py at <subcommand> [args...]
\end{verbatim}

\hypertarget{at-list}{%
\paragraph{\texorpdfstring{\texttt{at\ list}}{at list}}\label{at-list}}

List all available therapies with structural diagnosis data.

\begin{verbatim}
python3 rebis.py at list
\end{verbatim}

Output shows each therapy's name, structural distance (d) to the healthy
state, tier promotion (e.g.~O₀→O₂), and the primitives that need
promotion (Δ). Currently 10 therapies: schizophrenia, hiv, mrsa, mdd
(major depressive disorder), pcos, cf (cystic fibrosis),
gout\_elimination, gout\_combined, gout\_holistic, homeopathy.

\hypertarget{at-diagnose-disease}{%
\paragraph{\texorpdfstring{\texttt{at\ diagnose\ \textless{}disease\textgreater{}}}{at diagnose \textless disease\textgreater{}}}\label{at-diagnose-disease}}

Show full structural diagnosis --- the disease's 12-primitive tuple, its
distance from healthy/generic human, and the specific primitive
conflicts.

\begin{verbatim}
python3 rebis.py at diagnose schizophrenia
python3 rebis.py at diagnose mrsa
python3 rebis.py at diagnose cf
\end{verbatim}

Output includes: disease tuple, healthy tuple, per-primitive conflict
list, structural distance, and tier assessment.

\hypertarget{at-therapy-disease}{%
\paragraph{\texorpdfstring{\texttt{at\ therapy\ \textless{}disease\textgreater{}}}{at therapy \textless disease\textgreater{}}}\label{at-therapy-disease}}

Show the grammar-derived therapy protocol --- the exact promotion path
from disease state to healthy state.

\begin{verbatim}
python3 rebis.py at therapy schizophrenia
python3 rebis.py at therapy mrsa
python3 rebis.py at therapy hiv
\end{verbatim}

Output lists each primitive promotion step with the structural operation
required (e.g., φ̂: 𐑢→⊙, Ħ: 𐑓→𐑫), and compounds/molecules that effect
that promotion.

\hypertarget{at-compare-type_a-type_b}{%
\paragraph{\texorpdfstring{\texttt{at\ compare\ \textless{}type\_a\textgreater{}\ \textless{}type\_b\textgreater{}}}{at compare \textless type\_a\textgreater{} \textless type\_b\textgreater{}}}\label{at-compare-type_a-type_b}}

Compare two structural types --- show per-primitive alignment and
structural distance.

\begin{verbatim}
python3 rebis.py at compare schizophrenia mdd
python3 rebis.py at compare hiv mrsa
\end{verbatim}

\hypertarget{at-tensor-type_a-type_b}{%
\paragraph{\texorpdfstring{\texttt{at\ tensor\ \textless{}type\_a\textgreater{}\ \textless{}type\_b\textgreater{}}}{at tensor \textless type\_a\textgreater{} \textless type\_b\textgreater{}}}\label{at-tensor-type_a-type_b}}

Compute the tensor product of two structural types --- the composite's
absorbing structural type (⊙\_3 rule applies).

\begin{verbatim}
python3 rebis.py at tensor schizophrenia antidepressant
python3 rebis.py at tensor hiv antiviral
\end{verbatim}

\hypertarget{at-meet-type_a-type_b}{%
\paragraph{\texorpdfstring{\texttt{at\ meet\ \textless{}type\_a\textgreater{}\ \textless{}type\_b\textgreater{}}}{at meet \textless type\_a\textgreater{} \textless type\_b\textgreater{}}}\label{at-meet-type_a-type_b}}

Compute the meet (greatest lower bound) of two structural types --- the
shared structural floor.

\begin{verbatim}
python3 rebis.py at meet schizophrenia mdd
\end{verbatim}

\hypertarget{at-spectrum-type}{%
\paragraph{\texorpdfstring{\texttt{at\ spectrum\ \textless{}type\textgreater{}}}{at spectrum \textless type\textgreater{}}}\label{at-spectrum-type}}

Show the psychiatric spectrum mapping for a structural type --- how it
positions across mental health dimensions.

\begin{verbatim}
python3 rebis.py at spectrum schizophrenia
python3 rebis.py at spectrum mdd
\end{verbatim}

\hypertarget{at-operate-tuple-operation}{%
\paragraph{\texorpdfstring{\texttt{at\ operate\ \textless{}tuple\textgreater{}\ \textless{}operation\textgreater{}}}{at operate \textless tuple\textgreater{} \textless operation\textgreater{}}}\label{at-operate-tuple-operation}}

Apply a structural operation to a comma-separated 12-tuple. Operations
include: \texttt{calcination}, \texttt{dissolution},
\texttt{separation}, \texttt{conjunction}, \texttt{fermentation},
\texttt{distillation}, \texttt{coagulation}.

\begin{verbatim}
python3 rebis.py at operate --tuple "𐑼,𐑡,𐑩,𐑗,𐑱,𐑺,𐑲,𐑝,𐑢,𐑓,𐑙,𐑷" --operation calcination
\end{verbatim}

\hypertarget{scripts-1}{%
\subsubsection{\texorpdfstring{\texttt{scripts}}{scripts}}\label{scripts-1}}

Run standalone scripts that are not fully integrated as
\texttt{rebis.py} subcommands. These scripts live in \texttt{scripts/}
and are callable via this unified interface.

\begin{verbatim}
python3 rebis.py scripts <subcommand> [script_name] [args...]
\end{verbatim}

\hypertarget{scripts-list}{%
\paragraph{\texorpdfstring{\texttt{scripts\ list}}{scripts list}}\label{scripts-list}}

List all available scripts with line counts.

\begin{verbatim}
python3 rebis.py scripts list
\end{verbatim}

Shows each script name and its line count. Currently \textasciitilde21
scripts in \texttt{scripts/} including:

\begin{longtable}[]{@{}ll@{}}
\toprule
Script & Description\tabularnewline
\midrule
\endhead
\texttt{run\_antibody*} & Antibody CDR designer
(interactive)\tabularnewline
\texttt{mito\_pipeline*} & Mitochondrial genome pipeline
(v1--v3)\tabularnewline
\texttt{run\_msa} & Multiple sequence alignment\tabularnewline
\texttt{run\_pdb\_validation} & PDB structure validation\tabularnewline
\texttt{analyze\_validation} & Validation results
analyzer\tabularnewline
\texttt{gen\_univ\_map} & Diaschizic IUPAC generator\tabularnewline
\texttt{frobenius\_exact\_design*} & Frobenius exact design
tools\tabularnewline
\texttt{compute\_promotions} & Promotion path computer\tabularnewline
\texttt{psychedelic\_bridge*} & Psychedelic compound
bridge\tabularnewline
\texttt{serpentrod\_protein*} & SerpentRod standalone
scripts\tabularnewline
\texttt{ghost\_typer} & Demo segment runner for screen
recording\tabularnewline
\texttt{operations} & Alchemical operations\tabularnewline
\texttt{stress\_test\_proteins} & 34-test protein \& genetics stress
test\tabularnewline
\texttt{test\_genetics} & Full genetics test suite\tabularnewline
\texttt{\_demo\_*} & Demo scripts for each pillar\tabularnewline
\bottomrule
\end{longtable}

\hypertarget{scripts-run-name-args...}{%
\paragraph{\texorpdfstring{\texttt{scripts\ run\ \textless{}name\textgreater{}\ {[}args...{]}}}{scripts run \textless name\textgreater{} {[}args...{]}}}\label{scripts-run-name-args...}}

Run a script by name, optionally forwarding arguments.

\begin{verbatim}
python3 rebis.py scripts list                          # Show all scripts
python3 rebis.py scripts run run_antibody               # Interactive antibody designer
python3 rebis.py scripts run mito_pipeline              # MT gene pipeline
python3 rebis.py scripts run test_genetics              # Full B4 genetics suite
python3 rebis.py scripts run test_genetics --b4         # B4 lattice tests only
python3 rebis.py scripts run test_genetics --quick      # Quick test subset
python3 rebis.py scripts run stress_test_proteins       # 34-test stress suite
python3 rebis.py scripts run ghost_typer                # Demo typing (slow)
python3 rebis.py scripts run ghost_typer --fast         # Demo typing (fast)
python3 rebis.py scripts run ghost_typer --section 5    # Single section
python3 rebis.py scripts run operations                 # Alchemical operations
python3 rebis.py scripts run frobenius_exact_design     # Frobenius exact design
python3 rebis.py scripts run compute_promotions         # Promotion paths
python3 rebis.py scripts run run_msa                    # MSA analysis
python3 rebis.py scripts run run_pdb_validation         # PDB validation
python3 rebis.py scripts run analyze_validation         # Validation results
python3 rebis.py scripts run gen_univ_map              # IUPAC generator
python3 rebis.py scripts run psychedelic_bridge         # Psychedelic bridge

---

### `status` — Expanded

Shows all platform modules with file size, line count, and health check.

```bash
python3 rebis.py status
\end{verbatim}

Sample output (abbreviated):

\begin{verbatim}
╭──────────────────────┬───────┬─────────┬─────────────╮
│ Package              │ Files │    Size │ Root file   │
├──────────────────────┼───────┼─────────┼─────────────┤
│ ✅ ch3mpiler         │    19 │ 355,212 │ __init__.py │
│ ✅ clink             │    33 │ 685,758 │ __init__.py │
│ ✅ imas              │    11 │ 211,815 │ __init__.py │
│ ✅ materials         │    21 │ 551,018 │ __init__.py │
│ ✅ rhr_p4rky         │    32 │ 390,060 │ __init__.py │
│ ✅ serpentrod        │     3 │ 103,065 │ __init__.py │
│ ✅ shared            │     5 │  65,970 │ __init__.py │
│ ... (14 packages)    │       │         │             │
╰──────────────────────┴───────┴─────────┴─────────────╯
✅ shared/primitives.py: 13,467 bytes
✅ shared/IG_catalog.json: 2,050,279 bytes
14 packages + 21 scripts discovered
\end{verbatim}

Use this as your first command when setting up --- all packages should
show ✅. A ❌ means an import failure, usually from a missing dependency
(see Troubleshooting).

\hypertarget{verify-expanded}{%
\subsubsection{\texorpdfstring{\texttt{verify} ---
Expanded}{verify --- Expanded}}\label{verify-expanded}}

Imports every module and reports pass/fail. Includes Frobenius closure
checks where applicable.

\begin{verbatim}
python3 rebis.py verify
\end{verbatim}

This checks: - All 14 packages import cleanly - All \texttt{rhr\_p4rky/}
modules import (32 modules) - All shared utilities load - No circular
imports

Expected: 15+ lines all showing ✅. Run after any install, code change,
or \texttt{git\ pull}.

\begin{center}\rule{0.5\linewidth}{0.5pt}\end{center}

\hypertarget{comprehensive-run-target-reference}{%
\subsection{\texorpdfstring{Comprehensive \texttt{run} Target
Reference}{Comprehensive run Target Reference}}\label{comprehensive-run-target-reference}}

Discoverable via
\texttt{rebis.py\ run\ \textless{}target\textgreater{}}:

\begin{longtable}[]{@{}lll@{}}
\toprule
\begin{minipage}[b]{0.24\columnwidth}\raggedright
Target\strut
\end{minipage} & \begin{minipage}[b]{0.27\columnwidth}\raggedright
Command\strut
\end{minipage} & \begin{minipage}[b]{0.40\columnwidth}\raggedright
Description\strut
\end{minipage}\tabularnewline
\midrule
\endhead
\begin{minipage}[t]{0.24\columnwidth}\raggedright
\texttt{serpent\_rod}\strut
\end{minipage} & \begin{minipage}[t]{0.27\columnwidth}\raggedright
\texttt{rebis.py\ run\ serpent\_rod}\strut
\end{minipage} & \begin{minipage}[t]{0.40\columnwidth}\raggedright
V5 protein prediction --- test set\strut
\end{minipage}\tabularnewline
\begin{minipage}[t]{0.24\columnwidth}\raggedright
\texttt{ch3mpiler}\strut
\end{minipage} & \begin{minipage}[t]{0.27\columnwidth}\raggedright
\texttt{rebis.py\ run\ ch3mpiler\ -\/-smiles\ "\textless{}SMILES\textgreater{}"}\strut
\end{minipage} & \begin{minipage}[t]{0.40\columnwidth}\raggedright
Retrosynthetic compiler\strut
\end{minipage}\tabularnewline
\begin{minipage}[t]{0.24\columnwidth}\raggedright
\texttt{ch3mpiler\_serpentrod\_pipeline}\strut
\end{minipage} & \begin{minipage}[t]{0.27\columnwidth}\raggedright
\texttt{rebis.py\ run\ ch3mpiler\_serpentrod\_pipeline\ -\/-cas\ \textless{}num\textgreater{}}\strut
\end{minipage} & \begin{minipage}[t]{0.40\columnwidth}\raggedright
Catalytic site design\strut
\end{minipage}\tabularnewline
\begin{minipage}[t]{0.24\columnwidth}\raggedright
\texttt{antibody\_designer}\strut
\end{minipage} & \begin{minipage}[t]{0.27\columnwidth}\raggedright
\texttt{rebis.py\ run\ antibody\_designer}\strut
\end{minipage} & \begin{minipage}[t]{0.40\columnwidth}\raggedright
Advanced antibody design\strut
\end{minipage}\tabularnewline
\begin{minipage}[t]{0.24\columnwidth}\raggedright
\texttt{gene\_to\_protein\_pipeline}\strut
\end{minipage} & \begin{minipage}[t]{0.27\columnwidth}\raggedright
\texttt{rebis.py\ run\ gene\_to\_protein\_pipeline\ -\/-test}\strut
\end{minipage} & \begin{minipage}[t]{0.40\columnwidth}\raggedright
Full 7-stage pipeline\strut
\end{minipage}\tabularnewline
\begin{minipage}[t]{0.24\columnwidth}\raggedright
\texttt{psychedelic\_bridge}\strut
\end{minipage} & \begin{minipage}[t]{0.27\columnwidth}\raggedright
\texttt{rebis.py\ run\ psychedelic\_bridge}\strut
\end{minipage} & \begin{minipage}[t]{0.40\columnwidth}\raggedright
Psychedelic compound bridge\strut
\end{minipage}\tabularnewline
\begin{minipage}[t]{0.24\columnwidth}\raggedright
\texttt{diaschizic\_iupac}\strut
\end{minipage} & \begin{minipage}[t]{0.27\columnwidth}\raggedright
\texttt{rebis.py\ run\ diaschizic\_iupac}\strut
\end{minipage} & \begin{minipage}[t]{0.40\columnwidth}\raggedright
IUPAC generator\strut
\end{minipage}\tabularnewline
\begin{minipage}[t]{0.24\columnwidth}\raggedright
\texttt{decay\_chain}\strut
\end{minipage} & \begin{minipage}[t]{0.27\columnwidth}\raggedright
\texttt{rebis.py\ run\ decay\_chain}\strut
\end{minipage} & \begin{minipage}[t]{0.40\columnwidth}\raggedright
Nuclear decay IMASM analysis\strut
\end{minipage}\tabularnewline
\bottomrule
\end{longtable}

Use \texttt{rebis.py\ run\ list} to see the live list. Currently 8
discoverable targets.

\hypertarget{direct-run-scripts-not-discoverable-via-rebis.py-run}{%
\subsubsection{\texorpdfstring{Direct-run Scripts (not discoverable via
\texttt{rebis.py\ run})}{Direct-run Scripts (not discoverable via rebis.py run)}}\label{direct-run-scripts-not-discoverable-via-rebis.py-run}}

Use \texttt{python3\ scripts/\textless{}name\textgreater{}.py} or
\texttt{python3\ scripts/\textless{}name\textgreater{}.py\ {[}args{]}}:

\begin{longtable}[]{@{}llll@{}}
\toprule
\begin{minipage}[b]{0.17\columnwidth}\raggedright
Script\strut
\end{minipage} & \begin{minipage}[b]{0.24\columnwidth}\raggedright
Directory\strut
\end{minipage} & \begin{minipage}[b]{0.19\columnwidth}\raggedright
Command\strut
\end{minipage} & \begin{minipage}[b]{0.28\columnwidth}\raggedright
Description\strut
\end{minipage}\tabularnewline
\midrule
\endhead
\begin{minipage}[t]{0.17\columnwidth}\raggedright
\texttt{serpentrod}\strut
\end{minipage} & \begin{minipage}[t]{0.24\columnwidth}\raggedright
\texttt{serpentrod/}\strut
\end{minipage} & \begin{minipage}[t]{0.19\columnwidth}\raggedright
\texttt{python3\ -m\ serpentrod.protein\_v5}\strut
\end{minipage} & \begin{minipage}[t]{0.28\columnwidth}\raggedright
V5 protein prediction\strut
\end{minipage}\tabularnewline
\begin{minipage}[t]{0.17\columnwidth}\raggedright
\texttt{stress\_test\_proteins}\strut
\end{minipage} & \begin{minipage}[t]{0.24\columnwidth}\raggedright
\texttt{scripts/}\strut
\end{minipage} & \begin{minipage}[t]{0.19\columnwidth}\raggedright
\texttt{python3\ scripts/stress\_test\_proteins.py}\strut
\end{minipage} & \begin{minipage}[t]{0.28\columnwidth}\raggedright
34-test cross-pipeline stress test\strut
\end{minipage}\tabularnewline
\begin{minipage}[t]{0.17\columnwidth}\raggedright
\texttt{test\_genetics}\strut
\end{minipage} & \begin{minipage}[t]{0.24\columnwidth}\raggedright
\texttt{scripts/}\strut
\end{minipage} & \begin{minipage}[t]{0.19\columnwidth}\raggedright
\texttt{python3\ scripts/test\_genetics.py\ {[}-\/-b4/-\/-quick{]}}\strut
\end{minipage} & \begin{minipage}[t]{0.28\columnwidth}\raggedright
B4 lattice genetics test suite\strut
\end{minipage}\tabularnewline
\begin{minipage}[t]{0.17\columnwidth}\raggedright
\texttt{ghost\_typer}\strut
\end{minipage} & \begin{minipage}[t]{0.24\columnwidth}\raggedright
\texttt{scripts/}\strut
\end{minipage} & \begin{minipage}[t]{0.19\columnwidth}\raggedright
\texttt{python3\ scripts/ghost\_typer.py\ {[}-\/-fast/-\/-section\ N{]}}\strut
\end{minipage} & \begin{minipage}[t]{0.28\columnwidth}\raggedright
Demo typing for screen recording\strut
\end{minipage}\tabularnewline
\begin{minipage}[t]{0.17\columnwidth}\raggedright
\texttt{operations}\strut
\end{minipage} & \begin{minipage}[t]{0.24\columnwidth}\raggedright
\texttt{scripts/}\strut
\end{minipage} & \begin{minipage}[t]{0.19\columnwidth}\raggedright
\texttt{python3\ scripts/operations.py}\strut
\end{minipage} & \begin{minipage}[t]{0.28\columnwidth}\raggedright
Alchemical operations\strut
\end{minipage}\tabularnewline
\begin{minipage}[t]{0.17\columnwidth}\raggedright
\texttt{frobenius\_exact\_design}\strut
\end{minipage} & \begin{minipage}[t]{0.24\columnwidth}\raggedright
\texttt{scripts/}\strut
\end{minipage} & \begin{minipage}[t]{0.19\columnwidth}\raggedright
\texttt{python3\ scripts/frobenius\_exact\_design.py}\strut
\end{minipage} & \begin{minipage}[t]{0.28\columnwidth}\raggedright
Frobenius exact design\strut
\end{minipage}\tabularnewline
\begin{minipage}[t]{0.17\columnwidth}\raggedright
\texttt{compute\_promotions}\strut
\end{minipage} & \begin{minipage}[t]{0.24\columnwidth}\raggedright
\texttt{scripts/}\strut
\end{minipage} & \begin{minipage}[t]{0.19\columnwidth}\raggedright
\texttt{python3\ scripts/compute\_promotions.py}\strut
\end{minipage} & \begin{minipage}[t]{0.28\columnwidth}\raggedright
Promotion path computer\strut
\end{minipage}\tabularnewline
\begin{minipage}[t]{0.17\columnwidth}\raggedright
\texttt{run\_antibody}\strut
\end{minipage} & \begin{minipage}[t]{0.24\columnwidth}\raggedright
\texttt{scripts/}\strut
\end{minipage} & \begin{minipage}[t]{0.19\columnwidth}\raggedright
\texttt{python3\ scripts/run\_antibody.py}\strut
\end{minipage} & \begin{minipage}[t]{0.28\columnwidth}\raggedright
Antibody designer (interactive)\strut
\end{minipage}\tabularnewline
\begin{minipage}[t]{0.17\columnwidth}\raggedright
\texttt{run\_pdb\_validation}\strut
\end{minipage} & \begin{minipage}[t]{0.24\columnwidth}\raggedright
\texttt{scripts/}\strut
\end{minipage} & \begin{minipage}[t]{0.19\columnwidth}\raggedright
\texttt{python3\ scripts/run\_pdb\_validation.py}\strut
\end{minipage} & \begin{minipage}[t]{0.28\columnwidth}\raggedright
PDB structure validation\strut
\end{minipage}\tabularnewline
\begin{minipage}[t]{0.17\columnwidth}\raggedright
\texttt{analyze\_validation}\strut
\end{minipage} & \begin{minipage}[t]{0.24\columnwidth}\raggedright
\texttt{scripts/}\strut
\end{minipage} & \begin{minipage}[t]{0.19\columnwidth}\raggedright
\texttt{python3\ scripts/analyze\_validation.py}\strut
\end{minipage} & \begin{minipage}[t]{0.28\columnwidth}\raggedright
Validation results analyzer\strut
\end{minipage}\tabularnewline
\begin{minipage}[t]{0.17\columnwidth}\raggedright
\texttt{gen\_univ\_map}\strut
\end{minipage} & \begin{minipage}[t]{0.24\columnwidth}\raggedright
\texttt{scripts/}\strut
\end{minipage} & \begin{minipage}[t]{0.19\columnwidth}\raggedright
\texttt{python3\ scripts/gen\_univ\_map.py}\strut
\end{minipage} & \begin{minipage}[t]{0.28\columnwidth}\raggedright
Diaschizic IUPAC generator\strut
\end{minipage}\tabularnewline
\begin{minipage}[t]{0.17\columnwidth}\raggedright
\texttt{run\_msa}\strut
\end{minipage} & \begin{minipage}[t]{0.24\columnwidth}\raggedright
\texttt{scripts/}\strut
\end{minipage} & \begin{minipage}[t]{0.19\columnwidth}\raggedright
\texttt{python3\ scripts/run\_msa.py}\strut
\end{minipage} & \begin{minipage}[t]{0.28\columnwidth}\raggedright
Multiple sequence alignment\strut
\end{minipage}\tabularnewline
\begin{minipage}[t]{0.17\columnwidth}\raggedright
\texttt{mito\_pipeline{[}\_v2,\_v3{]}}\strut
\end{minipage} & \begin{minipage}[t]{0.24\columnwidth}\raggedright
\texttt{scripts/}\strut
\end{minipage} & \begin{minipage}[t]{0.19\columnwidth}\raggedright
\texttt{python3\ scripts/mito\_pipeline.py}\strut
\end{minipage} & \begin{minipage}[t]{0.28\columnwidth}\raggedright
Mitochondrial genome pipeline\strut
\end{minipage}\tabularnewline
\begin{minipage}[t]{0.17\columnwidth}\raggedright
\texttt{hadron\_belnap}\strut
\end{minipage} & \begin{minipage}[t]{0.24\columnwidth}\raggedright
\texttt{rhr\_p4rky/}\strut
\end{minipage} & \begin{minipage}[t]{0.19\columnwidth}\raggedright
\texttt{python3\ -m\ rhr\_p4rky.hadron\_belnap}\strut
\end{minipage} & \begin{minipage}[t]{0.28\columnwidth}\raggedright
Hadronic Belnap classification\strut
\end{minipage}\tabularnewline
\begin{minipage}[t]{0.17\columnwidth}\raggedright
\texttt{quark\_belnap}\strut
\end{minipage} & \begin{minipage}[t]{0.24\columnwidth}\raggedright
\texttt{rhr\_p4rky/}\strut
\end{minipage} & \begin{minipage}[t]{0.19\columnwidth}\raggedright
\texttt{python3\ -m\ rhr\_p4rky.quark\_belnap}\strut
\end{minipage} & \begin{minipage}[t]{0.28\columnwidth}\raggedright
Quark Belnap analysis\strut
\end{minipage}\tabularnewline
\begin{minipage}[t]{0.17\columnwidth}\raggedright
\texttt{gene}\strut
\end{minipage} & \begin{minipage}[t]{0.24\columnwidth}\raggedright
\texttt{gene\_imscriber/}\strut
\end{minipage} & \begin{minipage}[t]{0.19\columnwidth}\raggedright
\texttt{python3\ -m\ gene\_imscriber.engine}\strut
\end{minipage} & \begin{minipage}[t]{0.28\columnwidth}\raggedright
IG-native genetic compiler\strut
\end{minipage}\tabularnewline
\begin{minipage}[t]{0.17\columnwidth}\raggedright
\texttt{demo\_*}\strut
\end{minipage} & \begin{minipage}[t]{0.24\columnwidth}\raggedright
\texttt{demo\_scripts/}\strut
\end{minipage} & \begin{minipage}[t]{0.19\columnwidth}\raggedright
\texttt{python3\ demo\_scripts/\_demo\_\textless{}name\textgreater{}.py}\strut
\end{minipage} & \begin{minipage}[t]{0.28\columnwidth}\raggedright
Per-pillar demos\strut
\end{minipage}\tabularnewline
\bottomrule
\end{longtable}

\hypertarget{run-examples}{%
\subsubsection{Run Examples}\label{run-examples}}

Retrosynthesis with exact fragment SMILES:

\begin{verbatim}
python3 rebis.py run ch3mpiler --smiles "CC(=O)OC1=CC=CC=C1C(=O)O" --depth 3
\end{verbatim}

SerpentRod protein prediction:

\begin{verbatim}
python3 rebis.py run serpent_rod --seq "MALWMRLLPLLALLALWGPDPAAAFVNQHLCGSHLVEALYLVCGERGFFYTPKTR"
# Or directly:
python3 -m serpentrod.protein_v5
\end{verbatim}

Full genetics test suite:

\begin{verbatim}
python3 scripts/test_genetics.py --quick
python3 scripts/stress_test_proteins.py
\end{verbatim}

\begin{center}\rule{0.5\linewidth}{0.5pt}\end{center}

\hypertarget{alchemy-expanded-subcommand-reference}{%
\subsubsection{\texorpdfstring{\texttt{alchemy} --- Expanded Subcommand
Reference}{alchemy --- Expanded Subcommand Reference}}\label{alchemy-expanded-subcommand-reference}}

The alchemical bridge maps the 7 classical stages of the Great Work onto
12-primitive IG structural operations. 19 subcommands covering analysis,
design, synthesis, and retro-synthesis.

\begin{verbatim}
python3 rebis.py alchemy <subcommand> [options]
\end{verbatim}

\hypertarget{analysis}{%
\paragraph{Analysis}\label{analysis}}

\begin{longtable}[]{@{}lll@{}}
\toprule
\begin{minipage}[b]{0.32\columnwidth}\raggedright
Subcommand\strut
\end{minipage} & \begin{minipage}[b]{0.35\columnwidth}\raggedright
Description\strut
\end{minipage} & \begin{minipage}[b]{0.24\columnwidth}\raggedright
Example\strut
\end{minipage}\tabularnewline
\midrule
\endhead
\begin{minipage}[t]{0.32\columnwidth}\raggedright
\texttt{report}\strut
\end{minipage} & \begin{minipage}[t]{0.35\columnwidth}\raggedright
Full bridge status --- all available scroll treatises, tiers, molecular
analogs\strut
\end{minipage} & \begin{minipage}[t]{0.24\columnwidth}\raggedright
\texttt{rebis.py\ alchemy\ report}\strut
\end{minipage}\tabularnewline
\begin{minipage}[t]{0.32\columnwidth}\raggedright
\texttt{tier\ \textless{}name\textgreater{}}\strut
\end{minipage} & \begin{minipage}[t]{0.35\columnwidth}\raggedright
Analyze a specific tier --- structural type, molecular analogs,
operations\strut
\end{minipage} & \begin{minipage}[t]{0.24\columnwidth}\raggedright
\texttt{rebis.py\ alchemy\ tier\ O\_inf\_self\_modeling}\strut
\end{minipage}\tabularnewline
\begin{minipage}[t]{0.32\columnwidth}\raggedright
\texttt{scroll-family}\strut
\end{minipage} & \begin{minipage}[t]{0.35\columnwidth}\raggedright
List the scroll family --- treatises where φ̂=⊙ and Ω=ℤ\strut
\end{minipage} & \begin{minipage}[t]{0.24\columnwidth}\raggedright
\texttt{rebis.py\ alchemy\ scroll-family}\strut
\end{minipage}\tabularnewline
\begin{minipage}[t]{0.32\columnwidth}\raggedright
\texttt{suggest\ \textless{}tier\textgreater{}}\strut
\end{minipage} & \begin{minipage}[t]{0.35\columnwidth}\raggedright
Suggest molecular designs from a structural tier\strut
\end{minipage} & \begin{minipage}[t]{0.24\columnwidth}\raggedright
\texttt{rebis.py\ alchemy\ suggest\ O1\_pedagogical}\strut
\end{minipage}\tabularnewline
\begin{minipage}[t]{0.32\columnwidth}\raggedright
\texttt{trace\ \textless{}treatise\textgreater{}}\strut
\end{minipage} & \begin{minipage}[t]{0.35\columnwidth}\raggedright
Trace the Grand Sequence on a treatise's tuple --- step-by-step\strut
\end{minipage} & \begin{minipage}[t]{0.24\columnwidth}\raggedright
\texttt{rebis.py\ alchemy\ trace\ O\_inf\_self\_modeling}\strut
\end{minipage}\tabularnewline
\begin{minipage}[t]{0.32\columnwidth}\raggedright
\texttt{decode\ \textless{}term\textgreater{}}\strut
\end{minipage} & \begin{minipage}[t]{0.35\columnwidth}\raggedright
Decode a modern scientific term into its alchemical equivalent\strut
\end{minipage} & \begin{minipage}[t]{0.24\columnwidth}\raggedright
\texttt{rebis.py\ alchemy\ decode\ "electron\ microscope"}\strut
\end{minipage}\tabularnewline
\begin{minipage}[t]{0.32\columnwidth}\raggedright
\texttt{learn\ \textless{}modern\textgreater{}}\strut
\end{minipage} & \begin{minipage}[t]{0.35\columnwidth}\raggedright
Learn a modern term --- map it to the alchemical framework\strut
\end{minipage} & \begin{minipage}[t]{0.24\columnwidth}\raggedright
\texttt{rebis.py\ alchemy\ learn\ -\/-modern\ "catalyst"}\strut
\end{minipage}\tabularnewline
\bottomrule
\end{longtable}

\hypertarget{molecular-design}{%
\paragraph{Molecular Design}\label{molecular-design}}

\begin{longtable}[]{@{}lll@{}}
\toprule
\begin{minipage}[b]{0.32\columnwidth}\raggedright
Subcommand\strut
\end{minipage} & \begin{minipage}[b]{0.35\columnwidth}\raggedright
Description\strut
\end{minipage} & \begin{minipage}[b]{0.24\columnwidth}\raggedright
Example\strut
\end{minipage}\tabularnewline
\midrule
\endhead
\begin{minipage}[t]{0.32\columnwidth}\raggedright
\texttt{alchemical-mol\ \textless{}smiles\textgreater{}}\strut
\end{minipage} & \begin{minipage}[t]{0.35\columnwidth}\raggedright
Analyze a SMILES string through the alchemical bridge\strut
\end{minipage} & \begin{minipage}[t]{0.24\columnwidth}\raggedright
\texttt{rebis.py\ alchemy\ alchemical-mol\ "CN1C=NC2=C1C(=O)N(C(=O)N2C)C"}\strut
\end{minipage}\tabularnewline
\begin{minipage}[t]{0.32\columnwidth}\raggedright
\texttt{retro\ \textless{}smiles\textgreater{}}\strut
\end{minipage} & \begin{minipage}[t]{0.35\columnwidth}\raggedright
Alchemical retrosynthesis --- trace a molecule's formation path\strut
\end{minipage} & \begin{minipage}[t]{0.24\columnwidth}\raggedright
\texttt{rebis.py\ alchemy\ retro\ "CC(=O)OC1=CC=CC=C1C(=O)O"}\strut
\end{minipage}\tabularnewline
\begin{minipage}[t]{0.32\columnwidth}\raggedright
\texttt{host\ \textless{}smiles\textgreater{}}\strut
\end{minipage} & \begin{minipage}[t]{0.35\columnwidth}\raggedright
Host-guest analysis --- check if a molecule is a suitable host\strut
\end{minipage} & \begin{minipage}[t]{0.24\columnwidth}\raggedright
\texttt{rebis.py\ alchemy\ host\ "C1=CC=C(C=C1)C2=CC=CC=C2"}\strut
\end{minipage}\tabularnewline
\begin{minipage}[t]{0.32\columnwidth}\raggedright
\texttt{bind\ \textless{}smiles\textgreater{}}\strut
\end{minipage} & \begin{minipage}[t]{0.35\columnwidth}\raggedright
Guest binding analysis\strut
\end{minipage} & \begin{minipage}[t]{0.24\columnwidth}\raggedright
\texttt{rebis.py\ alchemy\ bind\ -\/-guest\ "CCO"}\strut
\end{minipage}\tabularnewline
\begin{minipage}[t]{0.32\columnwidth}\raggedright
\texttt{zosimos\ \textless{}name\textgreater{}}\strut
\end{minipage} & \begin{minipage}[t]{0.35\columnwidth}\raggedright
Zosimos procedure --- 7-stage alchemical synthesis path\strut
\end{minipage} & \begin{minipage}[t]{0.24\columnwidth}\raggedright
\texttt{rebis.py\ alchemy\ zosimos\ aspirin}\strut
\end{minipage}\tabularnewline
\bottomrule
\end{longtable}

\hypertarget{operations}{%
\paragraph{Operations}\label{operations}}

\begin{longtable}[]{@{}lll@{}}
\toprule
\begin{minipage}[b]{0.32\columnwidth}\raggedright
Subcommand\strut
\end{minipage} & \begin{minipage}[b]{0.35\columnwidth}\raggedright
Description\strut
\end{minipage} & \begin{minipage}[b]{0.24\columnwidth}\raggedright
Example\strut
\end{minipage}\tabularnewline
\midrule
\endhead
\begin{minipage}[t]{0.32\columnwidth}\raggedright
\texttt{operate}\strut
\end{minipage} & \begin{minipage}[t]{0.35\columnwidth}\raggedright
Apply an alchemical operation to a tuple\strut
\end{minipage} & \begin{minipage}[t]{0.24\columnwidth}\raggedright
\texttt{rebis.py\ alchemy\ operate\ -\/-tuple\ "𐑼,𐑡,𐑩,𐑗,𐑱,𐑺,..."\ -\/-operation\ calcination}\strut
\end{minipage}\tabularnewline
\begin{minipage}[t]{0.32\columnwidth}\raggedright
\texttt{greenfire}\strut
\end{minipage} & \begin{minipage}[t]{0.35\columnwidth}\raggedright
Analyze green fire potential of a substrate\strut
\end{minipage} & \begin{minipage}[t]{0.24\columnwidth}\raggedright
\texttt{rebis.py\ alchemy\ greenfire\ -\/-substrate\ "CCO"}\strut
\end{minipage}\tabularnewline
\begin{minipage}[t]{0.32\columnwidth}\raggedright
\texttt{wavelength}\strut
\end{minipage} & \begin{minipage}[t]{0.35\columnwidth}\raggedright
Compute alchemical wavelength of a treatise\strut
\end{minipage} & \begin{minipage}[t]{0.24\columnwidth}\raggedright
\texttt{rebis.py\ alchemy\ wavelength}\strut
\end{minipage}\tabularnewline
\begin{minipage}[t]{0.32\columnwidth}\raggedright
\texttt{portico}\strut
\end{minipage} & \begin{minipage}[t]{0.35\columnwidth}\raggedright
Portico analysis --- structural entry points\strut
\end{minipage} & \begin{minipage}[t]{0.24\columnwidth}\raggedright
\texttt{rebis.py\ alchemy\ portico}\strut
\end{minipage}\tabularnewline
\begin{minipage}[t]{0.32\columnwidth}\raggedright
\texttt{stilling}\strut
\end{minipage} & \begin{minipage}[t]{0.35\columnwidth}\raggedright
Stilling analysis --- equilibration pathways\strut
\end{minipage} & \begin{minipage}[t]{0.24\columnwidth}\raggedright
\texttt{rebis.py\ alchemy\ stilling}\strut
\end{minipage}\tabularnewline
\begin{minipage}[t]{0.32\columnwidth}\raggedright
\texttt{ladder}\strut
\end{minipage} & \begin{minipage}[t]{0.35\columnwidth}\raggedright
Ladder analysis --- tier climbing paths\strut
\end{minipage} & \begin{minipage}[t]{0.24\columnwidth}\raggedright
\texttt{rebis.py\ alchemy\ ladder}\strut
\end{minipage}\tabularnewline
\begin{minipage}[t]{0.32\columnwidth}\raggedright
\texttt{grand-sequence}\strut
\end{minipage} & \begin{minipage}[t]{0.35\columnwidth}\raggedright
Full grand sequence --- traverse all 7 operations\strut
\end{minipage} & \begin{minipage}[t]{0.24\columnwidth}\raggedright
\texttt{rebis.py\ alchemy\ grand-sequence}\strut
\end{minipage}\tabularnewline
\begin{minipage}[t]{0.32\columnwidth}\raggedright
\texttt{key\ \textless{}n\textgreater{}}\strut
\end{minipage} & \begin{minipage}[t]{0.35\columnwidth}\raggedright
Show key number n (1--12) --- each maps to a primitive\strut
\end{minipage} & \begin{minipage}[t]{0.24\columnwidth}\raggedright
\texttt{rebis.py\ alchemy\ key\ 9} (φ̂=⊙)\strut
\end{minipage}\tabularnewline
\begin{minipage}[t]{0.32\columnwidth}\raggedright
\texttt{opus}\strut
\end{minipage} & \begin{minipage}[t]{0.35\columnwidth}\raggedright
The Great Work summary --- show the full opus\strut
\end{minipage} & \begin{minipage}[t]{0.24\columnwidth}\raggedright
\texttt{rebis.py\ alchemy\ opus}\strut
\end{minipage}\tabularnewline
\bottomrule
\end{longtable}

\hypertarget{cdxml-expanded-subcommand-reference}{%
\subsubsection{\texorpdfstring{\texttt{cdxml} --- Expanded Subcommand
Reference}{cdxml --- Expanded Subcommand Reference}}\label{cdxml-expanded-subcommand-reference}}

Generate, verify, and clean CDXML (ChemDraw XML) files for molecules,
aptamers, and materials.

\begin{verbatim}
python3 rebis.py cdxml <subcommand> [options]
\end{verbatim}

\hypertarget{cdxml-generate}{%
\paragraph{\texorpdfstring{\texttt{cdxml\ generate}}{cdxml generate}}\label{cdxml-generate}}

Generate CDXML files from SMILES or predefined molecule names.

\begin{verbatim}
# Generate all CDXML in batch
python3 rebis.py cdxml generate --all

# Generate only small molecules
python3 rebis.py cdxml generate --molecules-only

# Single molecule from SMILES
python3 rebis.py cdxml generate --smiles "CN1C=NC2=C1C(=O)N(C(=O)N2C)C" --name caffeine

# Single molecule from predefined library
python3 rebis.py cdxml generate --molecule 5-nitro-bufotenin

# Custom annotation on canvas
python3 rebis.py cdxml generate --smiles "CCO" --name ethanol --annotation "C2H5OH - Solvent"

# Print CDXML to stdout
python3 rebis.py cdxml generate --smiles "CCO" --print
\end{verbatim}

\hypertarget{cdxml-verify}{%
\paragraph{\texorpdfstring{\texttt{cdxml\ verify}}{cdxml verify}}\label{cdxml-verify}}

Verify a CDXML file for structural correctness.

\begin{verbatim}
python3 rebis.py cdxml verify --file /path/to/molecule.cdxml
\end{verbatim}

\hypertarget{cdxml-clean}{%
\paragraph{\texorpdfstring{\texttt{cdxml\ clean}}{cdxml clean}}\label{cdxml-clean}}

Clean and reformat a CDXML file.

\begin{verbatim}
python3 rebis.py cdxml clean
\end{verbatim}

\hypertarget{updated-directory-structure-post-cleanup}{%
\subsection{Updated Directory Structure
(Post-Cleanup)}\label{updated-directory-structure-post-cleanup}}

After the v2.3.5 root cleanup, the top-level directory has been
organized into clearly-named subdirectories:

\begin{verbatim}
red-hot_rebis/
├── rebis.py               # CLI entry point (stays at root)
├── setup.py               # Package setup
├── Makefile               # Build/verify/test automation
├── README.md              # Project readme
├── MANUAL.md              # Manual page (man-style)
├── USER_GUIDE.md          # This guide — comprehensive usage
├── INDEX.md               # Static reference data
├── demo_scripts/          # Demo scripts (_demo_*.py, _real_demo.py)
├── scripts/               # Standalone scripts (ghost_typer, operations, stress_test, test_genetics)
├── fasta/                 # FASTA sequence files
├── data/                  # JSON data files
├── images/                # Images (PNG, SVG, diagrams)
├── _archive/              # Patents, commit_logs, old HTML cards
├── docs/                  # General documentation
├── rhr_p4rky/             # Paraconsistent kernel (32 modules)
├── shared/                # Shared primitives, catalog, utilities
├── serpentrod/            # Serpent's Rod — protein folding
├── ch3mpiler/             # CH3MPILER — retrosynthesis
├── pipeline/              # Auto-imscription pipeline
├── gene_imscriber/        # Gene Imscriber
├── clink/                 # CLINK Chain
├── imas/                  # IMASM compound pipeline
├── imasm_iterator/        # IMASM arrangement iterator
├── materials/             # Materials forge
├── biology/               # Biology simulations
├── therapeutics/          # Therapeutic design
├── alchemical_bridge/     # Alchemical operations
├── Ars_Therapeutica/      # Ars Therapeutica
├── natural_products/      # Natural products database
├── cdxml/                 # CDXML generation
├── popular_protein/       # Protein structure validation
├── pdb/                   # PDB reference files
├── designs/               # Design outputs
├── glossary/              # Glossary files
├── ig-docs/               # IG documentation
└── genetics_animations/   # SVG genetics visualizations
\end{verbatim}

All scripts (formerly root-level \texttt{.py} files) now live in
\texttt{scripts/} or \texttt{demo\_scripts/}. All data files live in
\texttt{data/}, all images in \texttt{images/}, all FASTA sequences in
\texttt{fasta/}. Patents and archival material are in
\texttt{\_archive/}.

\begin{center}\rule{0.5\linewidth}{0.5pt}\end{center}

\hypertarget{workflow-recipes}{%
\subsection{Workflow Recipes}\label{workflow-recipes}}

Complete end-to-end recipes for common tasks.

\hypertarget{recipe-1-design-a-novel-protein}{%
\subsubsection{Recipe 1: Design a Novel
Protein}\label{recipe-1-design-a-novel-protein}}

\begin{verbatim}
# 1. Verify the system is wired
python3 rebis.py status
python3 rebis.py verify

# 2. Check where proteins sit in the structural chain
python3 rebis.py clink layer 3  # Molecule layer
python3 rebis.py clink layer 4  # Cell layer (where proteins function)

# 3. Run SerpentRod
python3 rebis.py run serpent_rod

# 4. Or run with a custom sequence
python3 -m serpentrod.protein_v5 --seq "MALWMRLLPLLALLALWGPDPAAAFVNQHLCGSHLVEALYLVCGERGFFYTPKTR"

# 5. Validate against crystal structure
python3 scripts/popular_protein/comprehensive_comparison.py

# 6. Stress-test across all modules
python3 scripts/stress_test_proteins.py
\end{verbatim}

\hypertarget{recipe-2-discover-disconnections-for-a-molecule}{%
\subsubsection{Recipe 2: Discover Disconnections for a
Molecule}\label{recipe-2-discover-disconnections-for-a-molecule}}

\begin{verbatim}
# 1. Encode the molecule as an IMASM compound
python3 rebis.py imas compound --smiles "CC(=O)OC1=CC=CC=C1C(=O)O" --json

# 2. Run retrosynthesis
python3 rebis.py run ch3mpiler --smiles "CC(=O)OC1=CC=CC=C1C(=O)O" --depth 3

# 3. Find cross-domain structural neighbors
python3 rebis.py imas analogies --smiles "CC(=O)OC1=CC=CC=C1C(=O)O" --limit 10

# 4. Check the reaction mechanism
python3 rebis.py imas reaction --smiles "CC(=O)C1=CC=CC=C1" --product "CC(C1=CC=CC=C1)(C)O"
\end{verbatim}

\hypertarget{recipe-3-design-a-whole-organism}{%
\subsubsection{Recipe 3: Design a Whole
Organism}\label{recipe-3-design-a-whole-organism}}

\begin{verbatim}
# 1. View the CLINK chain from quarks to organism
python3 rebis.py clink layer 0
python3 rebis.py clink layer 4
python3 rebis.py clink layer 8

# 2. Check which tool bridges are available
python3 rebis.py pipeline bridges

# 3. Design from a starting layer
python3 rebis.py pipeline from-layer 3 8  # molecule → organism

# 4. Generate actionable outputs
python3 rebis.py pipeline actionable --organism mammal

# 5. Inspect the outputs
ls clink/datasets/organism_designs/organism_mammal_actionable/
\end{verbatim}

\hypertarget{recipe-4-crystal-guided-drug-discovery}{%
\subsubsection{Recipe 4: Crystal-Guided Drug
Discovery}\label{recipe-4-crystal-guided-drug-discovery}}

\begin{verbatim}
# 1. Start from a known compound
python3 rebis.py imas compound --smiles "CN(C)CCC1=CNC2=CC=CC=C12"

# 2. Find its crystal neighbors
python3 rebis.py imas analogies --smiles "CN(C)CCC1=CNC2=CC=CC=C12" --limit 10

# 3. Use the crystal designer to find ⊙-critical analogs
python3 scripts/frobenius_exact_design.py --smiles "CN(C)CCC1=CNC2=CC=CC=C12"

# 4. Register a promising candidate
python3 rebis.py imas register --smiles "CN(C)CCC1=CNC2=CC=C(O)C([N+](=O)[O-])=C12" --name "5-nitro-bufotenin"

# 5. Generate its CDXML for ChemDraw
python3 rebis.py cdxml generate --molecule 5-nitro-bufotenin
\end{verbatim}

\hypertarget{recipe-5-full-pipeline-smiles-to-organism}{%
\subsubsection{Recipe 5: Full Pipeline --- SMILES to
Organism}\label{recipe-5-full-pipeline-smiles-to-organism}}

\begin{verbatim}
# Takes a SMILES string through the entire pipeline:
# SMILES → IMASM → IG tuple → CLINK bridge → organism design

# Step 1: Encode the molecule
SMILES="CC(=O)OC1=CC=CC=C1C(=O)O"
python3 rebis.py imas compound --smiles "$SMILES"

# Step 2: Register it
python3 rebis.py imas register --smiles "$SMILES" --name aspirin

# Step 3: Design from molecule layer to organism
python3 rebis.py pipeline from-layer 3 8
\end{verbatim}

\hypertarget{recipe-6-structural-therapy-design}{%
\subsubsection{Recipe 6: Structural Therapy
Design}\label{recipe-6-structural-therapy-design}}

\begin{verbatim}
# 1. List available therapies
python3 rebis.py at list

# 2. Diagnose a condition
python3 rebis.py at diagnose mdd

# 3. Get the therapy protocol
python3 rebis.py at therapy mdd

# 4. Compare with related conditions
python3 rebis.py at compare mdd schizophrenia
\end{verbatim}

\hypertarget{recipe-7-materials-discovery}{%
\subsubsection{Recipe 7: Materials
Discovery}\label{recipe-7-materials-discovery}}

\begin{verbatim}
# 1. Forge all predefined materials
python3 rebis.py materials forge --all

# 2. Forge from an IMASM canonical
python3 rebis.py materials forge --name I_Dialetheic_Bootstrap

# 3. Run Frobenius metamaterial simulation
python3 rebis.py materials frobenius --size 30 --cycles 50

# 4. Run Ouroboric crack-healing simulation
python3 rebis.py materials ouroboric --grains 128 --stress 1000

# 5. Perform Sophick Mercury analysis
python3 rebis.py materials sophick --name eagle_9_sophick

# 6. Check Frobenius exactor pathways
python3 rebis.py materials exactor --name pathways
\end{verbatim}

\hypertarget{recipe-8-alchemical-analysis-of-a-molecule}{%
\subsubsection{Recipe 8: Alchemical Analysis of a
Molecule}\label{recipe-8-alchemical-analysis-of-a-molecule}}

\begin{verbatim}
# 1. Get the bridge report
python3 rebis.py alchemy report

# 2. Analyze a molecule alchemically
python3 rebis.py alchemy alchemical-mol "CN1C=NC2=C1C(=O)N(C(=O)N2C)C"

# 3. Trace its alchemical retrosynthesis
python3 rebis.py alchemy retro "CC(=O)OC1=CC=CC=C1C(=O)O"

# 4. Run the full Zosimos procedure
python3 rebis.py alchemy zosimos caffeine
\end{verbatim}

\hypertarget{recipe-9-full-platform-verification}{%
\subsubsection{Recipe 9: Full Platform
Verification}\label{recipe-9-full-platform-verification}}

\begin{verbatim}
# Run every verification the platform offers in sequence:

echo "=== STATUS ===" && python3 rebis.py status && \
echo "=== VERIFY ===" && python3 rebis.py verify && \
echo "=== CLINK ===" && python3 rebis.py clink layer 0 && \
echo "=== CLINK ===" && python3 rebis.py clink layer 8 && \
echo "=== PIPELINE BRIDGES ===" && python3 rebis.py pipeline bridges && \
echo "=== IMAS HUNT ===" && python3 rebis.py imas hunt --samples 10000 && \
echo "=== GENETICS ===" && python3 scripts/test_genetics.py --quick && \
echo "=== PROTEINS ===" && python3 rebis.py run serpent_rod && \
echo "=== MATERIALS ===" && python3 rebis.py materials forge --all && \
echo "=== ALL PASSED ==="
\end{verbatim}

\hypertarget{recipe-10-demo-for-screen-recording}{%
\subsubsection{Recipe 10: Demo for Screen
Recording}\label{recipe-10-demo-for-screen-recording}}

\begin{verbatim}
# Ghost Typer types everything at human pace
python3 scripts/ghost_typer.py        # Full demo (3 acts, ~15 min)
python3 scripts/ghost_typer.py --fast  # Fast mode (~3 min)
python3 scripts/ghost_typer.py --section 1  # Act I only
\end{verbatim}

\hypertarget{version-history}{%
\subsection{Version History}\label{version-history}}

\begin{longtable}[]{@{}lll@{}}
\toprule
\begin{minipage}[b]{0.34\columnwidth}\raggedright
Version\strut
\end{minipage} & \begin{minipage}[b]{0.23\columnwidth}\raggedright
Date\strut
\end{minipage} & \begin{minipage}[b]{0.34\columnwidth}\raggedright
Changes\strut
\end{minipage}\tabularnewline
\midrule
\endhead
\begin{minipage}[t]{0.34\columnwidth}\raggedright
v2.3.5\strut
\end{minipage} & \begin{minipage}[t]{0.23\columnwidth}\raggedright
2026-06-30\strut
\end{minipage} & \begin{minipage}[t]{0.34\columnwidth}\raggedright
\textbf{Root Cleanup + Manual Expansion Edition} --- Root directory
organized: loose files moved into \texttt{demo\_scripts/},
\texttt{scripts/}, \texttt{fasta/}, \texttt{data/}, \texttt{images/},
\texttt{\_archive/}; 0-byte \texttt{catalase} and Windows zone
identifier removed; USER\_GUIDE expanded with full \texttt{at}
(Ars\,Therapeutica) coverage (19 subcommands), expanded \texttt{scripts}
section, accurate \texttt{run} target reference, expanded
\texttt{alchemy} (19 subcommands) and \texttt{cdxml} references, updated
directory structure, 10 complete workflow recipes\strut
\end{minipage}\tabularnewline
\begin{minipage}[t]{0.34\columnwidth}\raggedright
v2.3.4\strut
\end{minipage} & \begin{minipage}[t]{0.23\columnwidth}\raggedright
2026-06-28\strut
\end{minipage} & \begin{minipage}[t]{0.34\columnwidth}\raggedright
\textbf{Rich Formatting + Exact SMILES Edition}\strut
\end{minipage}\tabularnewline
\begin{minipage}[t]{0.34\columnwidth}\raggedright
v2.3.3\strut
\end{minipage} & \begin{minipage}[t]{0.23\columnwidth}\raggedright
2026-06-27\strut
\end{minipage} & \begin{minipage}[t]{0.34\columnwidth}\raggedright
\textbf{Compound IMASM Edition}\strut
\end{minipage}\tabularnewline
\begin{minipage}[t]{0.34\columnwidth}\raggedright
v2.3.2\strut
\end{minipage} & \begin{minipage}[t]{0.23\columnwidth}\raggedright
2026-06-27\strut
\end{minipage} & \begin{minipage}[t]{0.34\columnwidth}\raggedright
rhr\_p4rky expansion: 32 modules\strut
\end{minipage}\tabularnewline
\begin{minipage}[t]{0.34\columnwidth}\raggedright
v2.3.1\strut
\end{minipage} & \begin{minipage}[t]{0.23\columnwidth}\raggedright
2026-06-24\strut
\end{minipage} & \begin{minipage}[t]{0.34\columnwidth}\raggedright
PDB format robustness\strut
\end{minipage}\tabularnewline
\begin{minipage}[t]{0.34\columnwidth}\raggedright
v2.3.0\strut
\end{minipage} & \begin{minipage}[t]{0.23\columnwidth}\raggedright
2026-06-21\strut
\end{minipage} & \begin{minipage}[t]{0.34\columnwidth}\raggedright
Ch3mpiler-SerpentRod v4 overhaul\strut
\end{minipage}\tabularnewline
\begin{minipage}[t]{0.34\columnwidth}\raggedright
v2.2.1\strut
\end{minipage} & \begin{minipage}[t]{0.23\columnwidth}\raggedright
2026-06-10\strut
\end{minipage} & \begin{minipage}[t]{0.34\columnwidth}\raggedright
Bug-fix release\strut
\end{minipage}\tabularnewline
\begin{minipage}[t]{0.34\columnwidth}\raggedright
v2.2\strut
\end{minipage} & \begin{minipage}[t]{0.23\columnwidth}\raggedright
2026-06-10\strut
\end{minipage} & \begin{minipage}[t]{0.34\columnwidth}\raggedright
Static-data commands moved to INDEX.md\strut
\end{minipage}\tabularnewline
\begin{minipage}[t]{0.34\columnwidth}\raggedright
v2.1\strut
\end{minipage} & \begin{minipage}[t]{0.23\columnwidth}\raggedright
2026-06-10\strut
\end{minipage} & \begin{minipage}[t]{0.34\columnwidth}\raggedright
Primitive names corrected\strut
\end{minipage}\tabularnewline
\begin{minipage}[t]{0.34\columnwidth}\raggedright
v2.0\strut
\end{minipage} & \begin{minipage}[t]{0.23\columnwidth}\raggedright
2026-05\strut
\end{minipage} & \begin{minipage}[t]{0.34\columnwidth}\raggedright
CLINK pipeline, materials forge\strut
\end{minipage}\tabularnewline
\begin{minipage}[t]{0.34\columnwidth}\raggedright
v1.0\strut
\end{minipage} & \begin{minipage}[t]{0.23\columnwidth}\raggedright
2026-04\strut
\end{minipage} & \begin{minipage}[t]{0.34\columnwidth}\raggedright
Initial release\strut
\end{minipage}\tabularnewline
\bottomrule
\end{longtable}

\begin{center}\rule{0.5\linewidth}{0.5pt}\end{center}

\emph{Guide maintained by Lando⊗⊙perator. Structural type of this
document: \(\langle 𐑼;𐑥;𐑾;𐑬;𐑐;𐑧;𐑔;𐑠;⊙;𐑖;𐑳;𐑴 \rangle\)}
